\documentclass[10pt,a4paper]{amsart}
\usepackage[margin=19mm]{geometry}
\usepackage{amsmath,amssymb,mathtools}
\usepackage[scr=rsfs,cal=euler]{mathalfa}
\usepackage{xfrac}
\usepackage[unicode,hidelinks,bookmarks=false]{hyperref}
\newcommand{\ie}{i.e.\ }
\newcommand{\cf}{cf.\ }
\newcommand{\resp}{respectively\ }
\newcommand{\Subsection}{\subsection}
\providecommand{\nobreakdash}{\nobreak\hbox{-}}
\setlength{\emergencystretch}{2em}
\multlinegap0pt
\usepackage{dsfont,mathrsfs,enumitem}
\newtheorem{theo}{Theorem}[section]
\newtheorem{lemma}[theo]{Lemma}
\newtheorem{prop}[theo]{Proposition}
\newtheorem{coro}[theo]{Corollary}
\theoremstyle{definition}
\newtheorem{defi}[theo]{Definition}
\newtheorem{rema}[theo]{Remark}
\newtheorem{exam}[theo]{Example}
\providecommand{\eg}{e.g.\ }
\def\N{\mathds N}
\def\Z{\mathds Z}
\def\Q{\mathds Q}
\def\R{\mathds R}
\def\C{\mathds C}
\def\epsilon{\varepsilon}
\def\phi{\varphi} 
\def\rho{\varrho}
\def\id{{\it id}}
\def\A{\mathds A}
\def\G{\mathds G}
\def\P{\mathds P}
\def\O{\mathcal O}

\newenvironment{equationi}{\par\noindent\vspace*{\abovedisplayskip}\begin{minipage}{\linewidth}
\begin{equation}}{\end{equation}\end{minipage}\hrule height0pt\leavevmode}


\DeclareMathOperator{\Mor}{Mor}
\DeclareMathOperator{\Hom}{Hom}
\DeclareMathOperator{\Aut}{Aut}
\DeclareMathOperator{\End}{End}
\DeclareMathOperator{\SHom}{{\mathscr H}\!\!\text{$om$}\,}
\DeclareMathOperator{\SEnd}{{\mathscr E}\!\text{$nd$}\,}
\DeclareMathOperator{\SRep}{{\mathscr R}\!\text{$ep$}\,}
\DeclareMathOperator{\UHom}{\underline{Hom}}
\DeclareMathOperator{\UEnd}{\underline{End}}
\DeclareMathOperator{\UAut}{\underline{Aut}}
\DeclareMathOperator{\coker}{coker}
\DeclareMathOperator{\St}{St}
\DeclareMathOperator{\GL}{GL}
\DeclareMathOperator{\PGL}{PGL}
\DeclareMathOperator{\Spec}{Spec}
\DeclareMathOperator{\Proj}{Proj}
\DeclareMathOperator{\conv}{conv}
\DeclareMathOperator{\interior}{int}
\DeclareMathOperator{\relint}{relint}
\DeclareMathOperator{\Pic}{Pic}
\DeclareMathOperator{\NS}{NS}
\DeclareMathOperator{\Lim}{Lim}

\newcommand*{\mk}{\mkern -1mu}
\newcommand*{\Mk}{\mkern -2mu}
\newcommand*{\mK}{\mkern 1mu}
\newcommand*{\MK}{\mkern 2mu}

\hypersetup{urlcolor=purple, linkcolor=blue, citecolor=red}


\newcommand*{\romanenumi}{\renewcommand*{\theenumi}{\roman{enumi}}}
\newcommand*{\Romanenumi}{\renewcommand*{\theenumi}{\Roman{enumi}}}
\newcommand*{\alphenumi}{\renewcommand*{\theenumi}{\alph{enumi}}}
\newcommand*{\Alphenumi}{\renewcommand*{\theenumi}{\Alph{enumi}}}

\title[Inverse-limit quiver moduli functor]{Quivers and moduli spaces of pointed curves of genus zero: original Proposition2.15}
\author{Mark Blume and Lutz Hille}
\date{Source: Algebraic Combinatorics4(1)(2021),89--124; DOI10.5802/alco.152}
\begin{document}\maketitle
\noindent\textit{Editorial scope.} The counted unit is the original Proposition2.15 described by the catalog\textquotesingle s locator2.15. The original arbitrary-base quiver-moduli definitions, local-ring, wall and cross-ratio constructions are retained as uncounted core. Genus-zero compactification and Hassett applications are omitted. All author mathematical expressions, hypotheses and complete proofs within the selected core are retained.
\subsection*{Notations and conventions}
We work over a base scheme $S$, usually we think of $S$ as $\Spec\Z$.
Products are over the base scheme $S$ unless otherwise specified. 
We use the theory of schemes as developed in~\cite{EGA, EGA1}. 
We sometimes write $y\in Y$ for a geometric point $y\colon\Spec k\to Y$.
We use the standard theory of algebraic stacks, see~\cite{SP}. 
We work with algebraic stacks in the fppf topology.
Concerning stability of quiver representations we use the sign convention 
opposite to~\cite{Ki94}. For $d,d'\in\Q^n$ we write $d'\leq d$ if and only if 
$d'_i\leq d_i$ for all $i$, and we write $d'<d$ if $d'\leq d$ and $d'\neq d$.



\section{Moduli spaces of quiver representations}
\label{sec:moduliquiver}

\subsection{Moduli of quiver representations over arbitrary schemes}

Let $Q$ be a finite connected quiver with set of vertices $Q_0$ and
set of arrows $Q_1$. 
For an arrow $\alpha\in Q_1$ we denote $s(\alpha)$ and $t(\alpha)$ the source
and target, \ie $s(\alpha)\stackrel{\alpha\:}{\longrightarrow}t(\alpha)$.

\begin{defi}
A {\it representation} of $Q$ of dimension vector $(d_q)_{q\in Q_0}$ over 
a scheme $Y$ is a tuple $V=((V_q)_{q\in Q_0},(\phi_\alpha)_{\alpha\in Q_1})$ 
where $V_q$ is a locally free $\O_Y$-module of rank $d_q$ and
$\phi_\alpha\colon V_{s(\alpha)}\to V_{t(\alpha)}$ a homomorphism of 
$\O_Y$-modules. A representation over $Y$ is called free if the underlying
$\O_Y$-module is free. The notions of homomorphisms of representations,
subrepresentations etc.\ are defined as usual.
\end{defi}

\begin{defi}
The {\it representation space} of $Q$ for the dimension vector $d=(d_q)_{q\in Q_0}$
over the base scheme $S$ is the affine space over $S$
\[
R(Q,d)=\prod_{\alpha\in Q_1}\UHom_S(\O_S^{\oplus d_{s(\alpha)}},
\O_S^{\oplus d_{t(\alpha)}})
\]
where $\UHom_S(\O_S^{\oplus d_{s(\alpha)}},\O_S^{\oplus d_{t(\alpha)}})$ 
is the Hom-scheme that represents the contravariant functor on $S$-schemes
$Y\mapsto\Hom_Y(\O_Y^{\oplus d_{s(\alpha)}},\O_Y^{\oplus d_{t(\alpha)}})$.
Equivalently, the scheme $R(Q,d)$ can be defined as the scheme that represents 
the contravariant functor on $S$-schemes
\[
Y\;\mapsto\;\big\{\;\,\text{free representations with fixed basis of $(Q,d)$ 
over $Y$}\,\big\}.
\]
On the representation space $R(Q,d)$ the group scheme over $S$
\[
G(d)=\prod_{q\in Q_0}\GL(d_q)
\]
operates by $(g_q)_q(\nu_\alpha)_\alpha=(g_{t(\alpha)}\nu_\alpha 
g_{s(\alpha)}^{-1})_\alpha$ where $(g_q)_q\in G(d)(Y)$,
$(\nu_\alpha)_\alpha\in R(Q,d)(Y)$ for $S$-schemes $Y$.
The diagonal subgroup $\,\G_m\cong\Delta\subseteq G(d)$ operates trivially, 
giving rise to a $\,\overline{G(d)}=G(d)/\Delta$-operation on $R(Q,d)$.
\end{defi}

By~\cite{Ki94} the following definition of (semi)stability is equivalent to 
(semi)stability in the sense of GIT of~\cite{MFK} (with $\prod_q\det(g_q)^{-\theta_q}$
the character corresponding to $\theta$).

\begin{defi}
For a representation $V$ of $Q$ of dimension vector $d=(d_q)_{q\in Q_0}$
over an algebraically closed field and a weight 
$\theta=(\theta_q)_{q\in Q_0}\in\Q^{Q_0}$, using the notations
$\dim(V)=(\dim(V_q))_{q\in Q_0}$ and $\theta(d)=\sum_{q\in Q_0}\theta_qd_q$, 
we define: %\\

\begin{enumerate}[label=(\alph*)]
\item \label{defi1.3_a}
%{\rm(a)} 
The representation $V$ is called {\it $\theta$-semistable} if
$\theta(d)=0$ and every subrepresentation $V'\subseteq V$ satisfies 
$\theta(\dim(V'))\leq 0$. %\\

\item \label{defi1.3_b}
%{\rm(b)} 
A $\theta$-semistable representation $V$ is called {\it $\theta$-stable} if 
the only subrepresentations $V'\subseteq V$ with $\theta(\dim(V'))=0$ are $V$ and $0$.
\end{enumerate}
\end{defi}

\begin{defi}
A representation $V$ of $Q$ of dimension vector $(d_q)_{q\in Q_0}$
over a scheme $Y$ is called {\it $\theta$-(semi)stable} if for all geometric points 
$y\colon\Spec k\to Y$ the representation $V_y=V\otimes_{\O_Y} k$ is 
$\theta$-(semi)stable.
\end{defi}

The functors on $S$-schemes 
\[
Y\;\mapsto\;\big\{\text{free $\theta$-(semi)stable representations with fixed 
basis of $(Q,d)$ over $Y$}\big\}
\]
are represented by open $G(d)$-invariant subschemes $R^\theta(Q,d)\!\subseteq\!R(Q,d)$
(parametrising $\theta$-semistable representations) and 
$R^{\theta\text{-}s}(Q,d)\!\subseteq\!R^\theta(Q,d)$ (parametrising $\theta$-stable
representations).

%\medskip

In order to construct moduli spaces of representations, following~\cite{Ki94} we 
consider GIT quotients by the (geometrically) reductive group scheme $G(d)$ or 
$\overline{G(d)}$ (see~\text{\cite[Section~9]{Al14}} for some results on 
(geometrically) reductive group schemes). %\linebreak
GIT over arbitrary locally noetherian schemes has been worked out first in~\cite{Ses77}, 
here we use the formulation of~\cite{Al14} in terms of adequate moduli spaces 
of Artin stacks.

%\medskip

Let $\mathcal M^\theta(Q,d)=R(Q,d)/_{\!\theta}\,G(d)$ 
the GIT quotient with respect to $\mathscr L=\O_{R(Q,d)}$ 
with the linearisation determined by the character $\theta$
(if $\theta$ is not integral replace $\theta$ by $l\theta$ for some $l\in\Z_{>0}$
such that $l\theta$ is integral), that is
\[\textstyle
\mathcal M^\theta(Q,d)
=\Proj_S\bigoplus_{k\geq 0}(p_*(\mathscr L^{\otimes k}))^{G(d)}
\]
where $p\colon R(Q,d)\to S$ is the structure morphism. We have the quotient morphism 
\[
R^\theta(Q,d)\to\mathcal X^\theta(Q,d)\stackrel{q\,}{\to}\mathcal M^\theta(Q,d).
\] 
where $\mathcal X^\theta(Q,d)=[R^\theta(Q,d)/\overline{G(d)}]$ is the 
quotient stack. The next theorem follows from the result in the affine case
\cite[Thm.~9.1.4]{Al14} as this property is local~\cite[Prop.~5.2.9]{Al14}.

\begin{theo}[\cite{Al14}]\label{thm:adequatemodulispace}
$q\colon\mathcal X^\theta(Q,d)\to\mathcal M^\theta(Q,d)$
is an adequate moduli space.
\end{theo}

In particular, we have the following properties (see~\cite{Al14}, \cf also~\cite{Ses77}):
\begin{enumerate}[label=(\arabic*)]
\item\label{aftertheo1.5_1} The homomorphism $\O_{\mathcal M^\theta(Q,d)}\to q_*\O_{\mathcal X^\theta(Q,d)}$
is an isomorphism.
\item\label{aftertheo1.5_2} $q$ is surjective, universally closed and universally submersive.
\item\label{aftertheo1.5_3} For an algebraically closed field $k$ the morphism $q$ induces a bijection
\[
[\mathcal X^\theta(Q,d)(k)]\,/\!\sim\;\;\to\;\;\mathcal M^\theta(Q,d)(k)
\]
where $\sim$ is the equivalence relation on the set of isomorphism classes
$[\mathcal X^\theta(Q,d)(k)]$ defined by $x_1\sim x_2$ if 
$\overline{\{x_1\}}\cap\overline{\{x_2\}}\neq\emptyset$ 
in $\mathcal X^\theta(Q,d)_s$ for $x_1,x_2\in\mathcal X^\theta(Q,d)(k)$ 
over $s\colon\Spec k\to S$.
\item\label{aftertheo1.5_4} Assume that $S$ is locally noetherian. Then $\mathcal M^\theta(Q,d)$ is of 
finite type over $S$ and $q_*\mathscr F$ is coherent for every coherent 
$\O_{\mathcal X^\theta(Q,d)}$-module $\mathscr F$.
\item\label{aftertheo1.5_5} $q$ is universal for morphisms to schemes, \ie for an $S$-scheme $Z$ 
the map
\[
\Mor_S(\mathcal M^\theta(Q,d),Z)\to\Mor_S(\mathcal X^\theta(Q,d),Z),
\quad f\mapsto f\circ q
\]
is a bijection.
\item\label{aftertheo1.5_6} Let $Y\to\mathcal M^\theta(Q,d)$ be a flat morphism of $S$-schemes.
Then $\mathcal X^\theta(Q,d)\times_{\mathcal M^\theta(Q,d)}Y\to Y$
is an adequate moduli space.
\end{enumerate}

Statement~\ref{aftertheo1.5_5} means that $R^\theta(Q,d)\to\mathcal M^\theta(Q,d)$ is a 
categorical quotient. Restricting to the open subscheme of stable points 
$R^{\theta\text{-}s}(Q,d)\subseteq R^\theta(Q,d)$ there is the following 
result corresponding to the notion of geometric quotient.

\begin{coro} 
Let $\mathcal X^{\theta\text{-}s}(Q,d)\!=\![R^{\theta\text{-}s}(Q,d)/\overline{G(d)}]$
and $\mathcal M^{\theta\text{-}s}(Q,d)$ the image of
$\mathcal X^{\theta\text{-}s}(Q,d)$ in $\mathcal M^\theta(Q,d)$. 
Then $\mathcal X^{\theta\text{-}s}(Q,d)\to\mathcal M^{\theta\text{-}s}(Q,d)$
is an adequate moduli space. For algebraically closed fields $k$ the map 
$[\mathcal X^{\theta\text{-}s}(Q,d)(k)]\to\mathcal M^{\theta\text{-}s}(Q,d)(k)$ 
is a bijection, where $[\mathcal X^{\theta\text{-}s}(Q,d)(k)]$ denotes the set of 
isomorphism classes of objects in $\mathcal X^{\theta\text{-}s}(Q,d)(k)$. 
\end{coro}

In addition to these standard GIT results, in our situation we have:

\begin{prop}\label{prop:Rs-Mtheta-torsor}
The operation of $\overline{G(d)}$ on $R^{\theta\text{-}s}(Q,d)$ is free
and the morphism
$R^{\theta\text{-}s}(Q,d)\to\mathcal M^{\theta\text{-}s}(Q,d)$ 
is a $\overline{G(d)}$-torsor in the fppf topology.
\end{prop}
\begin{proof}
We first show that the operation of $\overline{G(d)}$ has trivial stabilisers.
For an $S$-scheme $Y$ the stabiliser $\St(v)\subseteq G(d)_Y$ of an element 
$v\in R^{\theta\text{-}s}(Q,d)(Y)$ contains the diagonal group scheme 
$\Delta_Y\subseteq G(d)_Y$ and coincides with the automorphism group scheme 
$\UAut_Y(V)$ of the corresponding free representation $V=((V_q)_q,(\nu_\alpha)_\alpha)$ 
of $(Q,d)$ over $Y$.
The automorphism group scheme $\UAut_Y(V)$ is open dense in the endomorphism
scheme $\UEnd_Y(V)$. As the $\O_Y$-module $\SEnd_Y(V)$ is the kernel of 
the homomorphism of locally free sheaves
$\bigoplus_q\SEnd_Y(V_q)\to\bigoplus_\alpha\SHom_Y(V_{s(\alpha)},V_{t(\alpha)})$, 
$(\phi_q)_q\mapsto(\nu_\alpha\phi_{s(\alpha)}-\phi_{t(\alpha)}\nu_\alpha)_\alpha$,
we have $\UEnd_Y(V)=\A_Y(\mathscr C)$ with 
$\mathscr C=\coker\big(\bigoplus_\alpha\SHom_Y(V_{s(\alpha)},V_{t(\alpha)})^\vee
\!\to\bigoplus_q\SEnd_Y(V_q)^\vee\big)$.
Since $v$ is stable, all geometric fibres of $\St(v)=\UAut_Y(V)$ are of dimension $1$
(containing the diagonal group $\Delta$), therefore all geometric fibres 
of $\UEnd_Y(V)$ are $1$-dimensional affine spaces, so $\UEnd_Y(V)$ is an 
affine line bundle over $Y$.
It follows that $\St(v)=\UAut_Y(V)=\Delta_Y$.

We show that the operation is proper, that is 
$\Psi\colon\overline{G(d)}\times R^{\theta\text{-}s}(Q,d)\to 
R^{\theta\text{-}s}(Q,d)\times R^{\theta\text{-}s}(Q,d)$, $(g,v)\mapsto(gv,v)$ 
is a proper morphism, using the valuative criterion~\cite[II, (7.3.8)]{EGA}.
We can assume we have the base scheme $S=\Spec\Z$. Let $A$ be a discrete valuation
ring with maximal ideal $\mathfrak m$, field of fractions $K$ and $k=A/\mathfrak m$,
and let $t$ be a generator of $\mathfrak m$.
Given a morphism $\Spec A\to R^{\theta\text{-}s}(Q,d)\times R^{\theta\text{-}s}(Q,d)$ 
we have to show that the map
$\Mor_{R^{\theta\text{-}s}(Q,d)\times R^{\theta\text{-}s}(Q,d)}
(\Spec A,\overline{G(d)}\times R^{\theta\text{-}s}(Q,d))
\to\Mor_{R^{\theta\text{-}s}(Q,d)\times R^{\theta\text{-}s}(Q,d)}
(\Spec K,\overline{G(d)}\times R^{\theta\text{-}s}(Q,d))$
coming from the canonical homomorphism $A\to K$
is surjective (it is injective because the morphism is separated).
The morphism $\Spec A\to R^{\theta\text{-}s}(Q,d)\times R^{\theta\text{-}s}(Q,d)$ 
corresponds to a pair of representations $V'=((V'_q)_q,(\nu'_\alpha)_\alpha)$, 
$V=((V_q)_q,(\nu_\alpha)_\alpha)$ of $(Q,d)$ over $\Spec A$ and 
a morphism $\Spec K\to\overline{G(d)}\times R^{\theta\text{-}s}(Q,d)$ 
making the diagram arising from the above maps (the usual diagram when applying 
the valuative criterion) commutative corresponds to a pair $(g,V'')$ such that 
$V''=V_K$, $gV_K=V'_K$.
The element $g\in\overline{G(d)}(K)$ corresponds to a family of matrices 
with coefficients in $K$ up to a common factor in $K^*$,
so we can represent it by a family of matrices $(M_q)_q$ with coefficients 
in $A$ and choose bases of $V_q$ and $V'_q$ such that all $M_q$ are diagonal
with diagonal entries of the form $t^e$.
Let $l$ be the minimal integer such that $t^l$ occurs as a diagonal entry of some $M_q$.
Considering the matrices for the homomorphisms $\nu_\alpha,\nu'_\alpha$
with respect to the chosen bases, the $(i,j)$-entries for the two matrices for 
a given $\alpha$ differ by a factor $t^e\neq 1$ if the diagonal entries for the 
corresponding pairs of basis elements differ by such a factor. 
In this case modulo $\mathfrak m$ one of the two $(i,j)$-entries vanishes.
Let $W'_q\subseteq V'_q$ be the submodules generated by those basis elements 
which correspond to the diagonal entries of the form $t^l$, and similary $(U_q)_q$ 
the submodules of $(V_q)_q$ generated by the basis elements corresponding
to the remaining diagonal entries.
Then central fibres of the modules $W'_q\subseteq V'_q$ form a 
subrepresentation $W'_k\subseteq V'_k$ and the central fibres of the modules 
$U_q\subseteq V_q$ form a subrepresentation $U_k\subset V_k$. 
If these were proper nonzero subrepresentations, then, because $V_k$, $V'_k$ are 
$\theta$-stable, both subrepresentations would satisfy $\theta(W'_k)<0$, 
$\theta(U_k)<0$, which is impossible. 
Therefore $W'_k=V'_k$, $U_k=0$. It follows that the element $g\in\overline{G(d)}(K)$
can be represented by matrices with coefficients in $A\setminus\mathfrak m$.
These matrices are invertible as matrices over $A$ and form the required 
element in $\overline{G(d)}(A)$.

Being proper and quasifinite, the morphism $\Psi\colon\overline{G(d)}\times 
R^{\theta\text{-}s}(Q,d)\to R^{\theta\text{-}s}(Q,d)\times R^{\theta\text{-}s}(Q,d)$ 
is finite by~\cite[III (1), (4.4.2)]{EGA} or~\cite[IV (3), (8.11.1)]{EGA}.
The morphism $\Psi$ is a monomorphism because the operation has trivial stabilisers.
As finite epimorphisms of rings are surjective (see~\cite[Ch.~IV, Prop.~1.7]{La69}), 
it follows that $\Psi$ is a closed embedding, \ie the operation
$\overline{G(d)}\times R^{\theta\text{-}s}(Q,d)\to R^{\theta\text{-}s}(Q,d)$ is free.

The sheaf-quotient (fppf) $(R^{\theta\text{-}s}(Q,d)/\overline{G(d)})_{\text{fppf}}$,
\ie the quotient in the category of fppf-sheaves, is the sheaf associated 
to the presheaf $Y\mapsto R^{\theta\text{-}s}(Q,d)(Y)/\overline{G(d)}(Y)$ 
on the fppf-site of $S$-schemes (see \eg \cite[3.4.1]{Br14}). As the operation 
$\overline{G(d)}\times R^{\theta\text{-}s}(Q,d)\to R^{\theta\text{-}s}(Q,d)$ 
is free and a geometric quotient
$\pi\colon R^{\theta\text{-}s}(Q,d)\to R^{\theta\text{-}s}(Q,d)/\overline{G(d)}$ 
exists, by~\cite[Thm.~6]{An73} the canonical morphism
$(R^{\theta\text{-}s}(Q,d)/\overline{G(d)})_{\text{fppf}}\to 
R^{\theta\text{-}s}(Q,d)/\overline{G(d)}$ 
is an isomorphism and the morphism $\pi$ is fppf. 
Since the sheaf-quotient 
$(\mk R^{\theta\text{-}s}(\mk Q,\mk d)/\overline{G\mk (\mk d\mk )}\mk )_{\text{fppf}}$
has the property that $\overline{G(d)}\times R^{\theta\text{-}s}(Q,d)\to
R^{\theta\text{-}s}(Q,d)\times_{(R^{\theta\text{-}s}(Q,d)/
\overline{G(d)})_{\text{fppf}}}R^{\theta\text{-}s}(Q,d)$
is an isomorphism (see \eg \cite[3.4.5]{Br14}), it follows that the morphism 
of schemes $R^{\theta\text{-}s}(Q,d)\to R^{\theta\text{-}s}(Q,d)/\overline{G(d)}$ 
is a $\overline{G(d)}$-torsor in the fppf topology.
\end{proof}

The tautological $G(d)$-equivariant representation $U_{R^{\theta\text{-}s}(Q,d)}$
on $R^{\theta\text{-}s}(Q,d)$ can be made into a $\overline{G(d)}$-equivariant 
representation if and only if the dimension vector $d$ is indivisible, 
\ie the set of integers $d_q$ is coprime, see~\cite[Proof of Prop.~5.3]{Ki94}.
In this case by descent along the $\overline{G(d)}$-torsor
$R^{\theta\text{-}s}(Q,d)\to\mathcal M^{\theta\text{-}s}(Q,d)$ one obtains
a representation $U_{\mathcal M^{\theta\text{-}s}(Q,d)}$ 
on $\mathcal M^{\theta\text{-}s}(Q,d)$.

\begin{theo}\label{thm:modulireprquiver}
For an indivisible dimension vector $d=(d_q)_{q\in Q_0}$ the scheme 
$\mathcal M^{\theta\text{-}s}(Q,d)$ with universal representation 
$U_{\mathcal M^{\theta\text{-}s}(Q,d)}$ is a fine moduli space of $\theta$-stable 
representations of $(Q,d)$ in the sense that it represents the contravariant 
functor on $S$-schemes
\[
Y\;\mapsto\;
\big\{\text{$\,\theta$-stable representations of $(Q,d)$ over $Y\,$}\big\}\,/\!\sim
\]
of $\theta$-stable representations up to isomorphism over a Zariski cover.
\end{theo}
\begin{proof}
We denote the above functor $F^{\theta\text{-}s}(Q,d)$ and 
$q\colon R^{\theta\text{-}s}(Q,d)\to\mathcal M^{\theta\text{-}s}(Q,d)$
the quotient morphism.
We have the morphisms $\mathcal M^{\theta\text{-}s}(Q,d)\to 
F^{\theta\text{-}s}(Q,d)$, $(f\colon Y\to\mathcal M^{\theta\text{-}s}(Q,d))
\mapsto f^*U_{\mathcal M^{\theta\text{-}s}(Q,d)}$ 
and $F^{\theta\text{-}s}(Q,d)\to\mathcal M^{\theta\text{-}s}(Q,d)$,
$V\mapsto f_V$, where $f_V$ is locally defined as $f_V=q\circ\tilde{f}_V$ 
for free $V$ and $\tilde{f}_V\colon Y\to R^{\theta\text{-}s}(Q,d)$ is determined
by $V$ and a choice of basis of $V$, using that $R^{\theta\text{-}s}(Q,d)$ with
$U_{R^{\theta\text{-}s}(Q,d)}$ represents the functor of $\theta$-stable 
free representations with fixed bases.

To show that the composition 
$F^{\theta\text{-}s}(Q,d)\to\mathcal M^{\theta\text{-}s}(Q,d)
\to F^{\theta\text{-}s}(Q,d)$ is the identity, it is enough to show that 
$f_V^*U_{\mathcal M^{\theta\text{-}s}(Q,d)}\cong V$ for free representations $V$.
It is $f_V^*U_{\mathcal M^{\theta\text{-}s}(Q,d)}
=(q\circ \tilde{f}_V)^*U_{\mathcal M^{\theta\text{-}s}(Q,d)}
\cong\tilde{f}_V^*U_{R^{\theta\text{-}s}(Q,d)}\cong V$.

To show that the composition 
$\mathcal M^{\theta\text{-}s}(Q,d)\to F^{\theta\text{-}s}(Q,d)\to
\mathcal M^{\theta\text{-}s}(Q,d)$ is the identity,
it suffices to show that $f=f_{f^*U_{\mathcal M^{\theta\text{-}s}(Q,d)}}$
for $f\colon Y\to\mathcal M^{\theta\text{-}s}(Q,d)$ such that
$f^*U_{\mathcal M^{\theta\text{-}s}(Q,d)}$ is free.
Let $p_Y\colon Y'\to Y$ be an fppf morphism such that the $\overline{G(d)}$-torsor 
$R_Y=Y{}_f\!\!\times_q\,R^{\theta\text{-}s}(Q,d)\to Y$ becomes a trivial torsor 
$R_{Y'}\to Y'$ after base extension by $p_Y$. A section $Y'\to R_{Y'}$ gives a 
morphism $\tilde{f}'\colon Y'\to R^{\theta\text{-}s}(Q,d)$ which satisfies 
$q\circ\tilde{f}'=f\circ p_Y$.
We have $(\tilde{f}')^*U_{R^{\theta\text{-}s}(Q,d)}\cong
p_Y^*f^*U_{\mathcal M^{\theta\text{-}s}(Q,d)}\cong
p_Y^*(\tilde{f}_{f^*U_{\mathcal M^{\theta\text{-}s}(Q,d)}})^*
U_{R^{\theta\text{-}s}(Q,d)}$.
Therefore $\tilde{f}'$ coincides with 
$\tilde{f}_{f^*U_{\mathcal M^{\theta\text{-}s}(Q,d)}}\!\circ\,p_Y$
up to an element of $\overline{G(d)}$ that arises from comparison of
the two bases coming from $U_{R^{\theta\text{-}s}(Q,d)}$ of these isomorphic 
representations. Composing with $q$, it follows 
$f\circ p_Y=f_{f^*U_{\mathcal M^{\theta\text{-}s}(Q,d)}}\circ p_Y$,
and thus $f=f_{f^*U_{\mathcal M^{\theta\text{-}s}(Q,d)}}$ since 
$p_Y$ is an epimorphism.
\end{proof}

The equivalence relation on isomorphism classes of geometric points of 
$\mathcal X^\theta(Q,d)$ in~\ref{aftertheo1.5_3} after theorem~\ref{thm:adequatemodulispace},
or equivalently expressed as GIT-equivalence on $R^\theta(Q,d)(k)$
defined by $x_1\sim x_2$ if 
$\overline{G(d)_sx_1}\cap\overline{G(d)_sx_2}\neq\emptyset$
in $R^\theta(Q,d)_s$ for $x_1,x_2\in R^\theta(Q,d)(k)$ 
over $s\colon\Spec k\to S$, 
is by~\cite{Ki94} equivalent to S-equivalence of semistable representations 
as in the following definition.

\begin{defi}
Two $\theta$-semistable representations $V,V'$ of $Q$ over an algebraically 
closed field are called {\it S-equivalent} if there are filtrations in the 
category of $\theta$-semistable representations
$0=V_0\subsetneq\ldots\subsetneq V_l=V$ and $0=V'_0\subsetneq\ldots\subsetneq V'_l=V'$ 
with $\theta$-stable quotients such that 
$\bigoplus_iV_{i+1}/V_i\cong\bigoplus_iV'_{i+1}/V'_i$.
\end{defi}

In particular, $\theta$-stable representations are S-equivalent 
if and only if they are isomorphic. Note that the S in S-equivalence doesn't 
refer to the base scheme $S$.

\begin{coro}
Via the morphism $R^\theta(Q,d)\to\mathcal M^\theta(Q,d)$
any $\theta$-semistable representation of $(Q,d)$ over an $S$-scheme $Y$
determines a morphism $Y\to\mathcal M^\theta(Q,d)$ over $S$.
This way for a geometric point $s\colon\Spec k\to S$ the geometric points 
$\mathcal M^\theta(Q,d)_s(k)$ over $s$ are in bijection with 
the S-equivalence classes of $\theta$-semistable representations
of $(Q,d)$ over $k$.
\end{coro}


\subsection{Decomposition of the weight space and limits of moduli spaces of quiver representations}
\label{subsec:weightspace-limitfunctor}

In the case of GIT quotients of projective varieties $X$ over an algebraically closed 
field the chamber structure of the space $\NS^G(X)_\Q,$ where $\NS^G(X)$ is the 
N\'eron-Severi group of $G$-equivariant line bundles, has been described in 
\cite{Th96, DH98, Re00}, and in the case $X=R(Q,d)$ of quivers 
(without oriented cycles) in~\cite{HP02, Ch08}. 

%\medskip

Let $Q$ be a finite connected quiver, $d=(d_q)_{q\in Q_0}$ a dimension vector, 
$S$ be a scheme and $R(Q,d)$ the representation space over $S$. We have the space 
$\Q^{Q_0}$ of fractional linearisations of the line bundle $\O_{R(Q,d)}$. 

\begin{defi} We define the {\it weight space}
\[\textstyle
H(Q,d)=\big\{\theta\in\Q^{Q_0}\:|\:\theta(d)=0\big\}
\]
and
\[
C(Q,d)=\big\{\theta\in H(Q,d)\:|\:R^\theta(Q,d)\neq\emptyset\big\}.
\]
\end{defi}

For the finitely many dimension vectors $0<d'<d$ we define
\[
\begin{aligned}
H_{d'}&=\{\theta\in H(Q,d)\:|\:\theta(d')=0\},\\
H^{<0}_{d'}&=\{\theta\in H(Q,d)\:|\:\theta(d')<0\},\\
H^{\leq0}_{d'}&=\{\theta\in H(Q,d)\:|\:\theta(d')\leq0\}.\\
\end{aligned}
\]

\begin{rema}\label{rem:CQd}
Let the base scheme be $S=\Spec k$, $k$ an algebraically closed field.

\begin{enumerate}[label=(\arabic*)]
\item\label{rema1.12_1}
%(1) 
Existence of subrepresentations of a given dimension vector is a closed condition 
on $R(Q,d)$, see~\cite[Lemma~3.1]{Scho92}.
We call a dimension vector $0<d'<d$ a generic dimension vector if there is 
a non-empty open subscheme of $R(Q,d)$ of representations that have a 
subrepresentation of dimension vector $d'$, or equivalently, all representations
of dimension vector $d$ have a subrepresentation of dimension vector $d'$.
There is a dense open subset in $R(Q,d)$ of representations which only have
subrepresentations of generic dimension vectors. %\\

\item\label{rema1.12_2}
%(2) 
The set of weights that allow (semi)stable representations is determined by 
the generic dimension vectors:
the intersection $\bigcap_{\text{$d'\!\!$ generic}}H^{<0}_{d'}$ is the set of 
weights $\theta\in H(Q,d)$ such that $\theta$-stable representations exist, and
\[
C(Q,d)\;\;=\bigcap_{\text{$d'\!\!$ generic}}\!\!\!H^{\leq0}_{d'}
\] 
Thus $C(Q,d)$ is a convex polyhedral cone in $H(Q,d)$, and it is full-dimensional 
if there is a $\theta\in H(Q,d)$ such that $\theta$-stable representations exist.
If the quiver $Q$ has no oriented cycles then $C(Q,d)$ is strongly convex. %\\

\item\label{rema1.12_3}
%(3) 
By~\cite[Cor.~1]{Cr96}, improving the results of~\cite{Scho92} (see Thm.\ 3.3, 
Thm.\ 5.4 and the text following the proof of Thm.\ 5.4 in~\cite{Scho92}), the 
generic dimension vectors, and thus the cone $C(Q,d)$, are combinatorially determined. 
\end{enumerate}
\end{rema}

\begin{defi}
Weights $\theta,\theta'\in H(Q,d)$ are called {\it GIT equivalent} if 
\[R^\theta(Q,d)=R^{\theta'}(Q,d).\]
We denote the {\it GIT equivalence class} of $\theta$ by 
\[C_\theta=\{\theta'\in H(Q,d)\:|\:R^\theta(Q,d)=R^{\theta'}(Q,d)\}\]
and denote
\[\overline{C}_\theta=\{\theta'\in H(Q,d)\:|\:R^\theta(Q,d)\subseteq R^{\theta'}(Q,d)\}.\]
\end{defi}

%\medskip

For the standard results in the following proposition see~\cite{Ch08} (following
\cite{Re00}; proven for quivers without oriented cycles, but the results and 
methods of proof also apply to quivers with oriented cycles).

\begin{prop}\label{prop:GITclasses}
Let $S=\Spec k$, $k$ an algebraically closed field.
For $\theta\in C(Q,d)$ the set $\overline{C}_\theta$ is a convex polyhedral
cone and $\relint(\overline{C}_\theta)=C_\theta$. 
The cones $\overline{C}_\theta$ for $\theta\in C(Q,d)$ form a finite 
fan with support $C(Q,d)$.
\end{prop}

In the situation of proposition~\ref{prop:GITclasses} the decomposition of $C(Q,d)$ 
into GIT equivalence classes is determined by segments of the hyperplanes $H_{d'}$. 
The relevant part $W_{d'}\subseteq H_{d'}$ consists of those $\theta$ such that 
there is a $\theta$-semistable representation $V$ with a subrepresentation 
$V'\subset V$ of dimension vector $d'$ such that $\theta(V')=0$. Since in this case 
both $V'$ and $V/V'$ are also $\theta$-semistable, we have the equivalent condition 
that there is a $\theta$-semistable $V$ of the form $V\cong V'\oplus V''$ with 
$V',V''$ $\theta$-semistable of dimension vector $d',d-d'$, 
thus:
\[
W_{d'}=C(Q,d')\cap C(Q,d-d')
\]
Using the fact from remark~\ref{rem:CQd}.\ref{rema1.12_3} that the sets $C(Q,d)$ for any $d$ are 
combinatorially determined, this implies the following result.

\begin{prop}
The set $C(Q,d)\subseteq H(Q,d)$ and its decomposition into GIT equivalence classes 
over an algebraically closed field $k$ do not depend on the base field $k$.
\end{prop}

It follows that the above results concerning the decomposition of $H(Q,d)$, in particular 
proposition~\ref{prop:GITclasses}, also apply to the case of an arbitrary base scheme $S$. 
The weight space decomposition over $S$ can be determined by looking at an arbitrary fibre 
(or the situation over any algebraically closed field).

%\medskip

Assume that stable representations for some weight $\theta$ exist. Then $C(Q,d)$ is 
full-dimensional. The sets $W_{d'}$ of codimension one are called walls. 
Walls whose intersections with $\interior(C(Q,d))$ are nonempty are called inner walls, 
otherwise outer walls. A weight $\theta\in H(Q,d)$ is said to be generic if 
$R^\theta(Q,d)=R^{\theta\text{-}s}(Q,d)\neq\emptyset$. The GIT equivalence classes 
of generic weights are open sets called chambers, the connected components of weights
$\theta\in\interior C(Q,d)$ not contained in an inner wall. 

%\medskip

By variation of geometric invariant theory quotients (VGIT), \cf~\cite{Th96}, 
\cite{DH98}, quiver varieties of nearby weights are related. 
Let $C_{\theta_0},C_\theta\subset C(Q,d)$ be GIT equivalence 
classes such that $C_{\theta_0}\subseteq\overline{C}_\theta$.
Then the inclusion $R^\theta(Q,d)\subseteq R^{\theta_0}(Q,d)$
induces a morphism over $S$
\[\phi_{\theta,\theta_0}\colon\mathcal M^\theta(Q,d)\to\mathcal M^{\theta_0}(Q,d)\] 
The morphisms 
$\phi_{\theta,\theta_0}\colon\mathcal M^\theta(Q,d)\to\mathcal M^{\theta_0}(Q,d)$
for representatives of the GIT equivalence classes in $C(Q,d)$
form an inverse system and we can consider its inverse limit 
$\varprojlim_\theta\,\mathcal M^\theta(Q,d)$. 

\begin{prop}\label{prop:functorinvlim}
Let $Q$ be a quiver and $(d_q)_{q\in Q_0}$ be an indivisible dimension vector. 
Assume that stable representations for some weight $\theta$ exist. 
Then the functor $\varprojlim_\theta\,\mathcal M^\theta(Q,d)$ is isomorphic to the 
closed subfunctor of the product of $\mathcal M^\theta(Q,d)$ for generic $\theta$
\begin{equation}\label{eq:functorinvlim}
Y\;\mapsto\;
\left\{\left.
(\phi_\theta)_\theta\in\!\!\prod_{\theta \text{generic}}\!\!\!
\mathcal M^\theta(Q,d)(Y)
\;\;\right|\;
\begin{array}{l}
\forall\,\theta,\theta'\,\text{generic}\;\;\:
\forall\,\tau\in\overline{C}_\theta\cap\overline{C}_{\theta'}\colon\\
\phi_{\theta,\tau}\circ\phi_\theta=\phi_{\theta',\tau}\circ\phi_{\theta'}
\end{array}
\right\}
\end{equation}
\end{prop}
\begin{proof}
The limit functor is isomorphic to this functor since each $\tau\in C(Q,d)$ 
is contained in some $\overline{C}_\theta$ for a generic $\theta$.
This functor is a closed subfunctor of the product because $\mathcal M^\tau(Q,d)$
is a scheme separated over $S$ (use~\cite[(5.2.5)]{EGA1}).
\end{proof}



\section{Weight space decomposition and moduli spaces for some bipartite quivers}
\label{sec:weightspace-moduli-PnQn}

\subsection{The quivers \texorpdfstring{$P_n,Q_n$}{Pn,Qn} and the structure of the weight space}
\label{subsec:weightspacePnQn}

We define the quivers $P_n$ and $Q_n$ with fixed dimension vectors (\cf introduction).

\begin{defi}\ \\*[-1.2em]
\begin{enumerate}[label=(\alph*)]
\item\label{defi2.1_a} %{\rm (a)} 
Let
$Q_n=(\{p,q_1,\ldots,q_n\},\{\alpha_1,\ldots,\alpha_n\})$
where $s(\alpha_i)=q_i$, $t(\alpha_i)=p$ with dimension vector given by $d_p=2$ 
and $d_{q_i}=1$. %\\

\item\label{defi2.1_b} %{\rm (b)} 
Let $P_n=(\{p_1,p_2,q_1,\ldots,q_n\},\{\alpha_{1,1},\alpha_{1,2},\ldots,
\alpha_{n,1},\alpha_{n,2}\})$
where $s(\alpha_{i,j})=q_i$, $t(\alpha_{i,j})=p_j$ with dimension 
vector $(1,\ldots,1)$.
\end{enumerate}
\end{defi}

For these quivers with dimension vectors we determine the decomposition of the 
weight space as considered in subsection~\ref{subsec:weightspace-limitfunctor}. 
In the following when studying the structure of GIT classes in the cone 
$C(Q)\subset H(Q)$ for $Q=Q_n,P_n$ we restrict to the intersection with 
a certain affine hyperplane. In the case $Q_n$ this yields the hypersimplex 
$\Delta(2,n)$, in the case $P_n$ the product of simplices $\Delta^1\times\Delta^{n-1}$.

\begin{rema}\label{rem:Qn-weightspace} Decomposition of the weight space for $Q_n$. %\\

\begin{enumerate}[label=(\arabic*)]
\item\label{rema2.2_1}
%(1) 
A weight for the quiver $Q_n$ is given by a tuple 
$\theta=(\eta,\theta_1,\ldots,\theta_n)$ where %\linebreak 
$\theta_p=\eta$, $\theta_{q_i}=\theta_i$. 
The weight space is the $n$-dimensional hyperplane
$H(Q_n)= \{\theta\in\Q^{n+1}\:|\:2\eta=-\sum_i\theta_i\}\subset\Q^{n+1}$. 
Omitting the coordinate $\eta$ we have an isomorphism $H(Q_n)\cong\Q^n$, 
with basis $e_1,\ldots,e_n\in\Q^n$ and coordinates $\theta_1,\ldots,\theta_n$. 
A semistable representation exists if and only if 
$0\leq\theta_i\leq\frac{1}{2}\sum_{i=1}^n\theta_i$ for all $i$, 
this defines a full-dimensional strongly convex polyhedral cone $C(Q_n)\subset H(Q_n)$. 
The affine hyperplane $\{\sum_{i=1}^n\theta_i=2\}$ intersects 
this cone in the hypersimplex 
\[\hspace*{\leftmargini}\textstyle
C(Q_n)\cap\left\{\sum_{i=1}^n\theta_i=2\right\}=
\conv\big\{e_i+e_j\:|\:i\neq j\big\}\;=\;\Delta(2,n).
\]
\item\label{rema2.2_2} %(2) 
For $\theta=(\theta_1,\ldots,\theta_n)\in\Delta(2,n)$ a $\theta$-semistable 
representation that is not $\theta$-stable exists if and only if $\theta$ is 
contained in one of the following walls. The outer walls are of the form 
$W_i=\{\theta_i=0\}\cap\Delta(2,n)=\conv\{e_j+e_k\:|\:j,k\neq i\}
\cong\Delta(2,n-1)$ and 
$W_{\{\{i\},\{1,\ldots,n\}\setminus \{i\}\}}=\{\theta_i=1\}\cap\Delta(2,n)
=\conv\{e_i+e_j\:|\:j\in\{1\ldots,n\}\setminus\{i\}\}
\cong\Delta(1,n-1)=\Delta^{n-2}$.
The inner walls are 
\[\textstyle
\hspace*{\leftmargini}W_{\{J,J^\complement\}}\;=\;\big\{\theta\in\Delta(2,n)\:\big|\:\sum_{i\in J}\theta_i=1\big\}
\;=\;\conv\big\{e_i+e_j\:|\:i\in J,j\in J^\complement\big\}
\]
for $J\subset\{1,\ldots,n\}$, $2\leq|J|\leq n-2$, 
$J^\complement=\{1,\ldots,n\}\setminus J$. It is 
\[
\hspace*{\leftmargini}W_{\{J,J^\complement\}}\cong\Delta^{|J|-1}\times\Delta^{n-|J|-1}.
\] %\\

\item\label{rema2.2_3} %(3) 
An inner wall $W_{\{J,J^\complement\}}$ intersects the boundary of $\Delta(2,n)$
as follows:
$W_{\{J,J^\complement\}}\cap W_i=
\conv\{e_j+e_k\:|\:j\in J\setminus\{i\},k\in J^\complement\setminus\{i\}\}$ 
(\ie a wall $W_{\{J\setminus\{i\},J^\complement\setminus\{i\}\}}
\subset\Delta(2,\{1,\ldots,n\}\setminus\{i\})$)
and, assuming $i\in J$,
$W_{\{J,J^\complement\}}\cap W_{\{\{i\},\{i\}^\complement\}}
=\conv\{e_i+e_j\:|\:j\in J^\complement\}=\Delta^{|J^\complement|-1}$.%\\

\item\label{rema2.2_4} %(4) 
Distinct inner walls $W_{\{J_1,J_1^\complement\}},W_{\{J_2,J_2^\complement\}}
\subset\Delta(2,n)$ do not intersect in the interior of $\Delta(2,n)$ 
if and only if the two partitions of $\{1,\ldots,n\}$ into two subsets 
are compatible in the sense that there is an inclusion between the sets
$J_1,J_1^\complement$ and $J_2,J_2^\complement$.
In this case, if we assume $J_1\subset J_2$, the intersection
$W_{\{J_1,J_1^\complement\}}\cap W_{\{J_2,J_2^\complement\}}
=\conv\{e_j+e_k\:|\:j\in J_1,k\in J_2^\complement\}
=\Delta^{|J_1|-1}\times\Delta^{|J_2^\complement|-1}$
is a wall of the face $\Delta(2,J_1\cup J_2^\complement)$
of $\Delta(2,n)$.
\end{enumerate}
\end{rema}

We have the representation space $R(Q_n)=(\A^2)^n$ with the operation of 
$\GL(2)\times(\G_m)^n$. A free representation of $Q_n$ over a scheme $Y$ 
after a choice of a basis can be expressed as a tuple $V=(s_1,\ldots,s_n)$ 
of $n$ sections of $\A^2_Y\to Y$. 
We write $s_i\sim s_j$ if $s_i$ and $s_j$ are in the same $\G_m$-orbit.

\begin{lemma}\label{lemma:repr-pol-Qn}
Let $V=(s_1,\ldots,s_n)$ be a representation of $Q_n$ over an algebraically closed 
field that is semistable with respect to some weight in the interior of $\Delta(2,n)$,
in particular all $s_i$ are nonzero. Let the partition $\{1,\ldots,n\}=\bigsqcup_lJ_l$ 
be defined by $i,j\in J_l$ for some $l$ if and only if $s_i\sim s_j$.
The set $\Theta(V)=\{\theta\in\Delta(2,n)\:|\:\text{$V$ is $\theta$-semistable}\}$ 
is a closed convex polytope in $\Delta(2,n)$ with the following properties:

\begin{enumerate}[label=(\alph*)]
\item\label{lemma2.3_i} We have 
$\Theta(V)=\bigcap_l\{\theta\in\Delta(2,n)\:|\:\sum_{i\in J_l}\theta_i\leq 1\}$,
\ie $\Theta(V)$ is bounded within $\Delta(2,n)$ by the inner walls 
$W_{\{J_l,J_l^\complement\}}$ for $2\leq|J_l|\leq n-2$. Any two distinct of 
these walls do not intersect in the interior of $\Delta(2,n)$.

\item\label{lemma2.3_ii} The sets of vertices and edges of $\Theta(V)$ are subsets of the sets 
of vertices and edges of $\Delta(2,n)$.
$\Theta(V)=\conv\{e_i+e_j\:|\:\text{$i\in J_l$, $j\in J_{l'}$ for some $l\neq l'$}\}$.

\item\label{lemma2.3_iii} $\Theta(V)$ is full dimensional if and only if there are $\geq3$
classes $J_l$, or equivalently, $V$ is stable with respect to some
weight. Otherwise there are two classes $J$ and $J^\complement$ and 
$\Theta(V)=W_{\{J,J^\complement\}}$.
\end{enumerate}
\end{lemma}
\begin{proof} 
For~\ref{lemma2.3_i} note that the nontrivial subrepresentations of $V$ with one-dimensional 
subspace in $V_q$ are given by the sets $J_l$ (the spaces ($V_{q_j})_{j\in J_l}$ 
together with the one-dimensional subspace of $V_p$ determined by $(s_j)_{j\in J_l}$). 
That the walls do not intersect in the interior follows from remark 
\ref{rem:Qn-weightspace}.\ref{rema2.2_4}.
Concerning~\ref{lemma2.3_ii}, as the inner walls bounding $\Theta(V)$ do not intersect in the
interior of $\Delta(2,n)$, the first statement follows inductively using
remark~\ref{rem:Qn-weightspace}.\ref{rema2.2_3} and~\ref{rema2.2_4}. Then $\Theta(V)$ is the convex hull
of the vertices of $\Delta(2,n)$ it contains.
The statements in~\ref{lemma2.3_iii} are easy to show.
\end{proof}

The property that the sets of vertices and edges of $\Theta(V)$
are subsets of the sets of vertices and edges of $\Delta(2,n)$ means
that $\Theta(V)$ is a matroid polytope.
Matroid polytopes were defined in~\cite{GGMS87} (see also~\cite[(1.2.6)]{Ka93a}), 
and by~\cite[Thm.~4.1]{GGMS87} matroid polytopes in the hypersimplex $\Delta(k,n)$
correspond to matroids of rank $k$ on $\{1,\ldots,n\}$.
In the theory of GIT $\Theta(V)$ corresponds to the stability set of $V$ 
denoted $\Omega(V)$ in~\cite{Re00}.

%\medskip

Considering the quiver $P_n$ we have the representation space $R(P_n)=(\A^1\times\A^1)^n$ 
with the operation of $(\G_m\times\G_m)\times(\G_m)^n$. We write a free representation 
of $P_n$ over a scheme $Y$ after a choice of a basis as a tuple $V=(s_1,\ldots,s_n)$ 
of sections of $(\A^1\times\A^1)_Y\to Y$ where $s_i$ corresponds to 
$(\alpha_{i,1},\alpha_{i,2})$.
We set $s_i\sim s_j$ if $s_i$ and $s_j$ are in the same $\G_m$-orbit and add 
the sections $s_0=(0,1)$, $s_\infty=(1,0)$ which are only used to compare with 
sections $s_i$ up to $\sim$.

\begin{rema}\label{rem:Pn-weightspace} 
Decomposition of the weight space for $P_n$.%\\

\begin{enumerate}[label=(\arabic*)]
\item\label{rema2.4_1}
%(1) 
A weight for the quiver $P_n$ is an element
$\theta=(\eta_1,\eta_2,\theta_1,\ldots,\theta_n)$ where
$\theta_{p_i}=\eta_i$, $\theta_{q_i}=\theta_i$. We also write this as 
$\theta=\eta_1f_1+\eta_2f_2+\sum_i\theta_ie_i$ with 
$f_1,f_2,e_1,\ldots,e_n$ a basis of $\Q^{n+2}$.
The weight space is the $(n+1)$-dimensional hyperplane 
$H(P_n)=\{\sum_i\theta_i=-\eta_1-\eta_2\}\subset\Q^{n+2}$.
A semistable representation exists if and only if 
$\eta_1,\eta_2\leq0$ and $\theta_1,\ldots,\theta_n\geq 0$. 
This defines a full-dimensional cone $C(P_n)\subset H(P_n)$.
The affine hyperplane defined by $\eta_1+\eta_2=-1$ intersects this cone in
a product of simplices
\[
\hspace*{\leftmargini}
\begin{aligned}
C(P_n)\cap\{\eta_1\!+\!\eta_2\!=\!-1\}
&=\conv\{e_1\!-\!f_1,\ldots,e_n\!-\!f_1,e_1\!-\!f_2,\ldots,e_n\!-\!f_2\}\\
&=\Delta^1\times\Delta^{n-1}.
\end{aligned}
\]
\item\label{rema2.4_2}
%(2) 
The outer walls are of the form 
$W_{q_i}=\{\theta_i=0\}\cap(\Delta^1\times\Delta^{n-1})
\cong\Delta^1\times\Delta^{n-2}$ and 
$W_{p_j}=\{\eta_j=0\}\cap(\Delta^1\times\Delta^{n-1})\cong\Delta^{n-1}$.
The inner walls are 
\[\textstyle
\hspace*{\leftmargini}
\begin{aligned}
W_J&\textstyle =\left\{\theta\in\Delta^1\times\Delta^{n-1}\:\big|\:\sum_{i\in J}\theta_i=-\eta_1\right\}\\
&=\conv\big(\{e_i-f_1\:|\:i\in J\}\cup\{e_i-f_2\:|\:i\in J^\complement\}\big)
\end{aligned}
\]
for $J\subset\{1,\ldots,n\}$, $1\leq|J|\leq n-1$, $n\geq2$. 
It is $W_J\cong\Delta^{n-1}$.%\\
\item\label{rema2.4_3}
%(3) 
An inner wall $W_J$ intersects the boundary of $\Delta^1\times\Delta^{n-1}$
as follows:
$W_J\cap W_{q_i}=\conv\big(\{e_j-f_1\:|\:j\in J\setminus\{i\}\}\cup\{e_j-f_2\:|\:
j\in J^\complement\setminus\{i\}\}\big)$ 
(\ie a wall $W_{J\setminus\{i\}}\subset\Delta^1\times\Delta^{n-2}$)
and
$W_J\cap W_{q_1}=\conv\big(\{e_i-f_2\:|\:i\in J^\complement\}\big)
=\Delta^{|J^\complement|-1}$, 
$W_J\cap W_{q_2}=\conv\big(\{e_i-f_1\:|\:i\in J\}\big)
=\Delta^{|J|-1}$.%\\
\item\label{rema2.4_4}
%(4) 
Two inner walls $W_{J_1},W_{J_2}$ do not intersect in the interior of 
$\Delta^1\times\Delta^{n-1}$ if and only if $J_1\subsetneq J_2$
or $J_2\subsetneq J_1$.
\end{enumerate}
\end{rema}

\begin{rema}
The notation is such that for $\tau\in\relint(W_J)$ representations
$(s_1,\ldots,s_n)$ with $s_j\sim s_\infty=(1,0)$ for $j\in J$ and
$s_i\sim s_0=(0,1)$ for $i\in J^\complement$ are strictly %\linebreak 
\text{$\tau$-semistable.}
\end{rema}

The following lemma can be proven in the same way as lemma~\ref{lemma:repr-pol-Qn} 
using remark~\ref{rem:Pn-weightspace}.

\begin{lemma}\label{lemma:repr-pol-Pn}
Let $V\Mk = \Mk (s_1,\ldots,s_n)$ be a representation of $P_n$ over an algebraically closed 
field that is semistable with respect to some weight in the interior of %\linebreak
$\Delta^1\times\Delta^{n-1}$, in particular all $s_i$ are nonzero.
Let $J_0,J_\infty\subset\{1,\ldots,n\}$ be defined~by 
\[
J_0=\{i\:|\:s_i\sim s_0=(0,1)\},\ J_\infty=\{i\:|\:s_i\sim s_\infty=(1,0)\}.
\] The set 
$\Theta(V)=\{\theta\in\Delta^1\times\Delta^{n-1}\:|\:\text{$V$ is $\theta$-semistable}\}$
is a closed convex polytope in $\Delta^1\times\Delta^{n-1}$ with the following 
properties:
\begin{enumerate}[label=(\alph*)]
\item\label{lemma2.6_1} $\Theta(V)=\{\theta\in\Delta^1\times\Delta^{n-1}\:|\:
\sum_{i\in J_0}\theta_i\leq-\eta_2\}
\cap\{\theta\in\Delta^1\times\Delta^{n-1}\:|\:\sum_{i\in J_\infty}\theta_i\leq-\eta_1\}$,
\ie $\Theta(V)$ is bounded within $\Delta^1\times\Delta^{n-1}$ by the walls 
$W_{\!J_0^\complement},W_{J_\infty}$. These walls, if distinct, do not intersect in 
the interior of $\Delta^1\times\Delta^{n-1}$.
\item\label{lemma2.6_2} The set of vertices of $\Theta(V)$ is a subset of the set of vertices of 
$\Delta^1\times\Delta^{n-1}$. 
It is $\Theta(V)=
\conv\big(\{e_i-f_2\:|\:i\in J_0\}\cup\{e_i-f_1\:|\:i\in J_\infty\}\big)$.
\item\label{lemma2.6_3} $\Theta(V)$ is full dimensional if and only if 
$(J_0\cup J_\infty)^\complement\neq\emptyset$, or equivalently, $V$ 
is stable with respect to some weight. Otherwise 
$\Theta(V)=W_{\!J_0^\complement}=W_{J_\infty}$.
\end{enumerate}
\end{lemma}



\subsection{Relation between representations of \texorpdfstring{$Q_{n+2}$}{Qn+2} and \texorpdfstring{$P_n$}{Pn}}
\label{subsec:Qn+2-Pn}

We show that the structure of GIT equivalence classes in 
$\Delta^1\times\Delta^{n-1}\subset H(P_n)$ arises from the
structure of GIT equivalence classes in $\Delta(2,n+2)$ in the 
neighbourhood of some vertex $e_a+e_b$ and compare the corresponding
stacks and moduli spaces of representations.

%\medskip

The subset $\Delta(2,n+2)_{a,b}=\Delta(2,n+2)\cap
\{\theta_a+\theta_b\geq\sum_{i\neq a,b}\theta_i\}\setminus\{e_a+e_b\}$ 
is part of a cone with apex $e_a+e_b$ and we can project it onto
$W_{\{\{a,b\},\{a,b\}^\complement\}}=\Delta(2,n+2)\cap
\{\theta_a+\theta_b=\sum_{i\neq a,b}\theta_i\}$ by
\[
\begin{array}{rcl}
\Delta(2,n+2)_{a,b}&\to&W_{\{\{a,b\},\{a,b\}^\complement\}}\\
\theta=(\theta_a,\theta_b,(\theta_i)_{i\neq a,b})&\mapsto&
\bar{\theta}=(\lambda(\theta_a-1)+1,\lambda(\theta_b-1)+1,(\lambda\theta_i)_{i\neq a,b})\\
\end{array}
\]
where $\lambda=1/\sum_{i\neq a,b}\theta_i$.
The wall $W_{\{\{a,b\},\{a,b\}^\complement\}}=\conv(\{e_i+e_a\:|\:i\neq a,b\}\cup
\{e_i+e_b\:|\:i\neq a,b\})\subset\Delta(2,n+2)$ can be identified with 
$\Delta^1\times\Delta^{n-1}=\conv(\{e_i-f_2\:|\:i\neq a,b\}\cup
\{e_i-f_1\:|\:i\neq a,b\})\subset H(P_n)$ by 
\[
\begin{array}{rcl}
W_{\{\{a,b\},\{a,b\}^\complement\}}&\to&\Delta^1\times\Delta^{n-1}\subset H(P_n)\\[1mm]
\bar{\theta}=(\bar{\theta}_a,\bar{\theta}_b,(\bar{\theta}_i)_{i\neq a,b})&\mapsto&
\theta'=(\eta'_1=\bar{\theta}_a-1,\eta'_2=\bar{\theta}_b-1,
(\theta'_i=\bar{\theta}_i)_{i\neq a,b})\\
\end{array}
\]
We observe that the projections of the walls of $\Delta(2,n+2)$ near $e_a+e_b$
coincide with the walls of $\Delta^1\times\Delta^{n-1}\subset H(P_n)$. 
The walls meeting the interior of $\Delta(2,n+2)_{a,b}$ are the walls 
$W_{\{J,J^\complement\}}$ with $a\in J,b\in J^\complement$ or $a\in J^\complement,b\in J$. 
Their intersections with $W_{\{\{a,b\},\{a,b\}^\complement\}}$ are exactly the 
inner walls of $\Delta^1\times\Delta^{n-1}$: assume $a\in J, b\in J^\complement$, 
then the equation $\sum_{j\in J}\theta_j=1$ for $W_{\{J,J^\complement\}}$ 
is equivalent to the equation $\sum_{j\in J\setminus\{a\}}\theta'_j=-\eta'_1$
defining $W_{J\setminus\{a\}}\subset\Delta^1\times\Delta^{n-1}$.


\begin{prop}\label{prop:Qn+2-P_n}
Let $a,b\in\{1,\ldots,n+2\}$, $a\neq b$. %\\[1mm]

\begin{enumerate}[label=(\alph*)]
\item \label{prop2.7_a}
%{\rm (a)} 
The maps 
\begin{equationi}\label{eq:maprepr-Qn+2-P_n}
\begin{split}
\left\{
\begin{array}{l}
\text{free representations with fixed bases}\\
\text{$V=(s_1,\ldots,s_{n+2})$ of $Q_{n+2}$ over $Y$}\\
\text{such that $s_a\cap 0=s_b\cap 0=\emptyset$}\\ 
\text{and $\G_ms_a\!\cap\G_ms_b=\emptyset$}
\end{array}\right\}%\\
%\hspace*{20mm}
%\begin{aligned}
&\rightarrow
\left\{
\begin{array}{l}
\text{free representations}\\
\text{with fixed bases of $P_n$}\\
\text{over $Y$}
\end{array}\right\}
\\%[9mm]
V=(s_1,\ldots,s_{n+2})&\mapsto V'=(s'_i)_{i\neq a,b}
%\end{aligned}
\end{split}
\end{equationi}

\noindent
for schemes $Y$ over a base scheme $S$, where $s'_i$ arises from $s_i$ by the 
base change that transforms $s_a$ to $(1,0)$ and $s_b$ to $(0,1)$, induce an 
isomorphism of stacks
\begin{equationi}\label{eq:isom-Qn+2-Pn}
\begin{aligned}[t]
&[R(Q_{n+2})_{a,b}/((\GL(2)\times(\G_m)^n)/\G_m)]\\
&\hspace*{35mm}\;\to\;
[R(P_n)/((\G_m\times\G_m)\times(\G_m)^n)/\G_m)]
\end{aligned}
\end{equationi}%%%ICI

\noindent
where $R(Q_{n+2})_{a,b}=R(Q_{n+2})\setminus(\{s_a=0\}\cup\{s_b=0\}\cup
\{\G_m s_a=\G_ms_b\})$. %\\[1mm]



\item\label{prop2.7_b} %{\rm (b)} 
We have the map
\begin{equationi}\label{eq:proj-Qn+2-Pn}
\begin{aligned}[t]
\Delta(2,n+2)_{a,b}&\to\Delta^1\times\Delta^{n-1}\subset H(P_n)\\
\theta&\mapsto\theta'\Mk =\Mk (\eta'_1\Mk =\Mk \lambda(\theta_a\Mk -\Mk 1),\eta'_2\Mk =\Mk \lambda(\theta_b\Mk -\Mk 1),
(\theta'_i\Mk =\Mk \lambda\theta_i)_{i\Mk \neq \Mk a,b})
\end{aligned}
\end{equationi}

\noindent
where $\lambda=1/\sum_{i\neq a,b}\theta_i$.
On the fibres of this map the $\theta$-(semi)stable locus in
$R(Q_{n+2})_{a,b}$ does not change. %\\

\item\label{prop2.7_c}
%{\rm (c)} 
Let $\theta\in W_{\{\{a,b\},\{a,b\}^\complement\}}\subset\Delta(2,n+2)$ and 
$\theta'\in\Delta^1\times\Delta^{n-1}\subset H(P_n)$ its image under the
map~\eqref{eq:proj-Qn+2-Pn}.
Let $l\in\N_{>0}$ such that $l\theta$ and $l\theta'$ are integral. Then
the equivariant line bundles $\O_{R(Q_{n+2})_{a,b}}$ with $l\theta$-linearisation 
and $\O_{R(P_n)}$ with $l\theta'$-linearisation define isomorphic line bundles 
on the quotient stacks under the identification~\eqref{eq:isom-Qn+2-Pn}.
\end{enumerate}
\end{prop}


\begin{proof}
\ref{prop2.7_a}~The maps on the sets of morphisms are given by
composition with the base change maps that transform $s_a,s_b$ to $(1,0),(0,1)$
and forgetting the automorphisms of the spaces corresponding to the vertices $q_a,q_b$.
The inverse functor is given on the objects by maps $(s_i)_{i\neq a,b}\mapsto(s_i)_i$
setting $s_a=(1,0)$, $s_b=(0,1)$ and on the morphisms by taking those
automorphisms of the spaces corresponding to the vertices $q_a,q_b$ such that
$s_a,s_b$ remain fixed. One checks that both compositions of these
two functors are isomorphic to the identity functors.

\ref{prop2.7_b}~The fibres are the line segments $[\bar{\theta},e_a+e_b]\setminus\{e_a+e_b\}$
for $\bar{\theta}\in W_{\{\{a,b\},\{a,b\}^\complement\}}$.
All $\theta$ contained in this line segment are elements of the same GIT equivalence class 
except $\bar{\theta}$. But $R^{\bar{\theta}}(Q_{n+2})$ differs from 
$R^\theta(Q_{n+2})$ for the other $\theta$ only outside $R(Q_{n+2})_{a,b}$.

\ref{prop2.7_c}~The isomorphism~\eqref{eq:isom-Qn+2-Pn} identifies the structure sheaves of the 
two stacks. We have to verify that the additional multiplications by the corresponding 
characters coincide. It suffices to consider the groupoids of $Y$-valued points for 
$S$-schemes $Y$ of the two given groupoid schemes. The elements of 
$((\GL(2)\times(\G_m)^n)/\G_m)(Y)$ fixing $s_a=(1,0)$, $s_b=(0,1)$ are of the 
form $\Big(\left(\begin{smallmatrix}\alpha_a&0\\0&\alpha_b\end{smallmatrix}\right),
\alpha_a,\alpha_b,(\alpha_i)_{i\neq a,b}\Big)$,
the corresponding element
$(\beta_1,\beta_2,(\alpha_i)_{i\neq a,b})\in((\G_m\times\G_m)\times(\G_m)^n)/\G_m)(Y)$ 
satisfies $\beta_1=\alpha_a$, $\beta_2=\alpha_b$.
Applying the character corresponding to 
$l\theta=(l\eta,l\theta_1,\ldots,l\theta_{n+2})$ we obtain
$(\alpha_a\alpha_b)^{l\eta}\alpha_a^{l\theta_a}\alpha_b^{l\theta_b}
\prod_{i\neq a,b}\alpha_i^{l\theta_i}$ and applying $l\theta'$ we have 
$\beta_1^{l\eta'_1}\beta_2^{l\eta'_2}\prod_{i\neq a,b}\alpha_i^{l\theta'_i}$.
As $\eta=-1$, this coincides for $\theta'_i=\theta_i$, $\eta'_1=\theta_a-1$,
$\eta'_2=\theta_b-1$.
\end{proof}

%\pagebreak

\begin{coro}\label{cor:Qn+2-P_n}
Let $a,b\in\{1,\ldots,n+2\}$, $a\neq b$.%\\

\begin{enumerate}[label=(\alph*)]
%{\rm (a)} 
\item\label{coro2.7_a}
For $\theta\in\Delta(2,n+2)_{a,b}\setminus W_{\{\{a,b\},\{a,b\}^\complement\}}$
a free representation $V=(s_1,\ldots,s_{n+2})$ of $Q_{n+2}$ is $\theta$-(semi)stable 
if and only if $V$ satisfies $s_a\cap 0=\emptyset$, $s_b\cap 0=\emptyset$, 
$\G_m s_a\cap\G_ms_b=\emptyset$ and its image under~\eqref{eq:maprepr-Qn+2-P_n} 
is a $\theta'$-(semi)stable representation of $P_n$, where $\theta'$ is the image of 
$\theta$ under the map~\eqref{eq:proj-Qn+2-Pn}. %\\

\item\label{coro2.7_b}
%{\rm (b)} 
The map~\eqref{eq:proj-Qn+2-Pn} defines a bijection between
the GIT equivalence classes in $\Delta^1\times\Delta^{n-1}$ and in
$\Delta(2,n+2)_{a,b}\setminus W_{\{\{a,b\},\{a,b\}^\complement\}}$ 
(resp.\ $ W_{\{\{a,b\},\{a,b\}^\complement\}}$). %\\

\item\label{coro2.7_c}
%{\rm (c)} 
For $\theta\in\Delta(2,n+2)_{a,b}\setminus W_{\{\{a,b\},\{a,b\}^\complement\}}$
the isomorphism~\eqref{eq:isom-Qn+2-Pn} induces an isomorphism of stacks
\[
\hspace*{\leftmargini}\begin{aligned}
[R^\theta(Q_{n+2})/((\GL(2)&\times(\G_m)^n)/\G_m)]\\
&\;\to\;
[R^{\theta'}(P_n)/((\G_m\times\G_m)\times(\G_m)^n)/\G_m)]\hspace*{\leftmargini}
\end{aligned}
 \]
and thus an isomorphism of their moduli spaces
\[
\hspace*{\leftmargini}
\mathcal M^\theta(Q_{n+2})\;\to\;\mathcal M^{\theta'}(P_n).
\]
These isomorphisms determine an isomorphism of the inverse systems of the moduli spaces 
$\mathcal M^\theta(Q_{n+2})$ for $\theta\in\Delta(2,n+2)_{a,b}\setminus 
W_{\{\{a,b\},\{a,b\}^\complement\}}$ and $\mathcal M^{\theta'}(P_n)$ 
for $\theta'\in\Delta^1\times\Delta^{n-1}$. 
\end{enumerate}
\end{coro}




\subsection{Moduli spaces for \texorpdfstring{$P_n,Q_n$}{Pn,Qn} and GIT quotients of products of \texorpdfstring{$\P^1$}{P1}'s}
\label{subsec:quivermoduli-GITP1}

We compare the moduli spaces of representations of $Q_n$ and $P_n$ to GIT quotients 
of $(\P^1)^n$ by the quotient groups 
$((\GL(2)\times(\G_m)^n)/\G_m)/((\G_m\times(\G_m)^n)/\G_m)\cong\PGL(2)$ and
$(((\G_m)^2\times(\G_m)^n)/\G_m)/((\G_m\times(\G_m)^n)/\G_m)\cong\G_m$.

%\medskip

We have $\NS((\P^1)^n)=\Pic((\P^1)^n)\cong\Z^n$, where 
$(\theta_1,\ldots,\theta_n)\in\Z^n$ corresponds to the line bundle 
$\O_{(\P^1)^n}(\theta_1,\ldots,\theta_n)=\O_{\P^1}(\theta_1)\boxtimes\ldots\boxtimes
\O_{\P^1}(\theta_n)$. A line bundle 
$\O_{(\P^1)^n}(\theta_1,\ldots,\theta_n)$ with 
$\theta_i\in 2\Z$ for all $i$ has a $\PGL(2)$-linearisation ($\O_{\P^1}(-2)$
being the canonical sheaf of $\P^1$), and if an element of $\Pic((\P^1)^n)$ has a 
$\PGL(2)$-linearisation then this linearisation is unique. Thus we can identify 
$\NS^{\PGL(2)}((\P^1)^n)_\Q=\Pic^{\PGL(2)}((\P^1)^n)_\Q\cong\Q^n$ 
with $H(Q_n)$ by 
$(\theta_1,\ldots,\theta_n)\leftrightarrow(\eta=-\frac{1}{2}\sum_i\theta_i,\theta_1,
\ldots,\theta_n)$.

\begin{prop}\label{prop:Q_n-GITP1}
Let $R^*(Q_n)=R(Q_n)\setminus\bigcup_i\{s_i=0\}\cong(\A^2\setminus\{0\})^n$. %\\



\begin{enumerate}[label=(\alph*)]
\item\label{prop2.9_a}
%{\rm (a)} 
The morphism $R^*(Q_n)\to(\P^1)^n$ induces an isomorphism of stacks
\begin{equationi}\label{eq:isom-Q_n-GITP1}
[R^*(Q_n)/((\GL(2)\times(\G_m)^n)/\G_m)]\;\to\;[(\P^1)^n/\PGL(2)].
\end{equationi}

\item\label{prop2.9_b}
%{\rm (b)} 
Let $\theta \Mk = \Mk (\eta,\mk \theta_{\mk 1},\mk \ldots,\mk \theta_{\mk n}\mk )\Mk \! \in\Mk \!H\mk (\mk Q_{\mk n}\mk )\mk $ be integral. 
Then the line bundle %\break 
$\mk \O_{\mk (\mk \P^{\mk 1}\mk )^{\mk n}}\mk (\mk \theta_{\mk 1}\mk ,\mk \ldots\mk ,\mk \theta_{\mk n}\mk )$ has a 
(unique) $\PGL(2)$-linearisation and the equivariant line bundles $\O_{R^*(Q_n)}$ with 
$\theta$-linearisation and $\O_{(\P^1)^n}(\theta_1,\ldots,\theta_n)$ define 
isomorphic line bundles on the quotient stacks under the identification 
in~\eqref{eq:isom-Q_n-GITP1}.
\end{enumerate}
\end{prop}
\begin{proof}
Consider the operation of the subgroup $(\G_m)^n\cong(\G_m\times(\G_m)^n)/\G_m
\subset(\GL(2)\times(\G_m)^n)/\G_m$ on $R^*(Q_n)\cong(\A^2\setminus\{0\})^n$.
The operation of $\G_m$ on $\A^2\setminus\{0\}$ is free and $\A^2\setminus\{0\}\to\P^1$
a $\G_m$-torsor. 
The structure sheaf on $\A^2\setminus\{0\}$ with $\G_m$-linearisation by $\theta$ 
descends to the line bundle $\O_{\P^1}(\theta)$. Thus $R^*(Q_n)\to(\P^1)^n$ is a 
$(\G_m)^n$-torsor and $\O_{R^*(Q_n)}$ with $(\G_m)^n$-linearisation by 
$(\theta_1,\ldots,\theta_n)$ descends to the line bundle 
$\O_{(\P^1)^n}(\theta_1,\ldots,\theta_n)$. %\\

\ref{prop2.9_a}~A $Y$-valued point of $[R^*(Q_n)/((\GL(2)\times(\G_m)^n)/\G_m)]$ for some 
$S$-scheme $Y$ corresponds to a $(\GL(2)\times(\G_m)^n)/\G_m$-torsor
$E\to Y$ with an equivariant morphism $E\to R^*(Q_n)$.
Taking quotients by $(\G_m)^n\subset (\GL(2)\times(\G_m)^n)/\G_m$ 
we obtain a $\PGL(2)$-equivariant morphism $E/(\G_m)^n\to(\P^1)^n$
where $E/(\G_m)^n\to Y$ is a $\PGL(2)$-torsor. The diagram
\[
\begin{array}{ccc}
E&\longrightarrow&(\A^2\setminus\{0\})^n\\
\downarrow&&\downarrow\\
E/(\G_m)^n&\longrightarrow&(\P^1)^n\\
\end{array}
\]
is cartesian and its vertical arrows are $(\G_m)^n$-torsors. For a $Y$-valued point of 
$[(\P^1)^n/\PGL(2)]$, a $\PGL(2)$-torsor $\overline{E}\to Y$ with equivariant morphism
to $(\P^1)^n$, we construct a $(\GL(2)\times(\G_m)^n)/\G_m$-torsor over $Y$ with an 
equivariant morphism to $(\A^2\setminus\{0\})^n$ by taking the fibred product 
$\overline{E}\times_{(\P^1)^n}(\A^2\setminus\{0\})^n$.
One verifies that this way two functors between
$[R^*(Q_n)/((\GL(2)\times(\G_m)^n)/\G_m)]$ and $[(\P^1)^n/\PGL(2)]$
are defined whose compositions are isomorphic to the identity functors.

\ref{prop2.9_b}~As the categories of sheaves do not change under stackification we may work with 
the prestack $[R^*(Q_n)/((\GL(2)\times(\G_m)^n)/\G_m)]^\text{pre}$ of trivial 
$(\GL(2)\times(\G_m)^n)/\G_m$-torsors $E\to Y$ with an equivariant morphism 
$E\to R^*(Q_n)$, and its image in $[(\P^1)^n/\PGL(2)]$ under the isomorphism~\eqref{eq:isom-Q_n-GITP1}. 
Objects in $[R^*(Q_n)/((\GL(2)\times(\G_m)^n)/\G_m)]^\text{pre}$ over an 
$S$-scheme $Y$ correspond to morphisms $Y\to R^*(Q_n)$, a morphism 
$(g,y)\colon Y\to (\GL(2)\times(\G_m)^n)/\G_m\times R^*(Q_n)$ gives an arrow 
$y\to gy$ over $\id_Y$.

Let $\mathscr L$ be the line bundle on 
$[R^*(Q_n)/((\GL(2)\times(\G_m)^n)/\G_m)]^\text{pre}$ defined by the 
equivariant line bundle $\O_{R^*(Q_n)}$ with linearisation by $\theta$, in 
particular $\mathscr L(y)=\Gamma(Y,y^*\O_{R^*(Q_n)})$ for $y\colon Y\to R^*(Q_n)$. 
We describe the push-forward $\overline{\mathscr L}$ with respect to the isomorphism~\eqref{eq:isom-Q_n-GITP1}, 
which by definition is given by 
$\overline{\mathscr L}(\overline{y})
=\varprojlim_{\overline{y'}\to\overline{y}}\mathscr L(y')$ where $\overline{y'}$ is the 
image of the object $y'$ under~\eqref{eq:isom-Q_n-GITP1}.
For open embeddings $\overline{y}\colon\overline{U}\hookrightarrow(\P^1)^n$ and 
$y\colon U\hookrightarrow R^*(Q_n)$ such that $U\subseteq R^*(Q_n)$ is the preimage of
$\overline{U}\subseteq(\P^1)^n$ we have
$\overline{\mathscr L}(\overline{y})
=\varprojlim_{\overline{y'}\to\overline{y}}\mathscr L(y')
=\varprojlim_{h\in(\G_m)^n(U),\overline{hy}\to\overline{y}}\mathscr L(hy)
=\Gamma(U,\O_{R^*(Q_n)})^{(\G_m)^n}
=\Gamma(\overline{U},\O_{(\P^1)^n}(\theta_1,\ldots,\theta_n))$.
Further we have the natural restriction maps. Thus the line bundle 
$\overline{\mathscr L}$ comes from the line bundle 
$\O_{(\P^1)^n}(\theta_1,\ldots,\theta_n)$ on $(\P^1)^n$ with some $\PGL(2)$-linearisation.
\end{proof}

\begin{coro}\ \\*[-1.2em]\label{cor:reprQn-P1n-GIT}

\begin{enumerate}[label=(\alph*)]
\item\label{coro2.10_a}
%{\rm (a)} 
We have an identification
$\Pic^{\PGL(2)}((\P^1)^n)\Mk \cong \Mk H(Q_n)\cap\Z^{n+1}$ by 
$(\theta_1,\ldots,\theta_n)\leftrightarrow(\eta=-\frac{1}{2}\sum_i\theta_i,\theta_1,
\ldots,\theta_n)$.

\item\label{coro2.10_b}
%{\rm (b)} 
For $(\eta,\theta_1,\ldots,\theta_n)\in H(Q_n)$ such that $\theta_i>0$ 
for all $i$ a free representation $V=(s_1,\ldots,s_n)$ of $Q_n$ is 
$(\eta,\theta_1,\ldots,\theta_n)$-(semi)stable if and only if 
$s_i\cap 0=\emptyset$ for all $i$ and its image under~\eqref{eq:isom-Q_n-GITP1} 
is a $(\theta_1,\ldots,\theta_n)$-(semi)stable 
section of $(\P^1)^n$.

\item\label{coro2.10_c}
%{\rm (c)} 
The image of $\interior C(Q_n)\Mk \subset \Mk H(Q_n)$ under the isomorphism 
$\Pic^{\PGL(2)}((\P^1)^n)_\Q\Mk \cong H(Q_n)$ is the $\PGL(2)$-ample cone in 
$\Pic^{\PGL(2)}((\P^1)^n)_\Q$. The GIT classes in these open cones coincide.

\item\label{coro2.10_d}
%{\rm (d)} 
For $\theta\in H(Q_n)\cong\Pic^{\PGL(2)}((\P^1)^n)_\Q$ such that 
$\theta_i>0$ for all $i$ the isomorphism~\eqref{eq:isom-Q_n-GITP1} induces 
isomorphisms of the stacks of $\theta$-(semi)stable points, of their moduli spaces 
and their inverse systems.
\end{enumerate}
\end{coro}


One can also compare these stacks with $[G(2,n)/((\G_m)^n/\G_m)]$,
\cf~\cite[2.4]{Ka93a}.

%\medskip

Similarly, one can show the analogous results for the quiver $P_n$.
In the case of the quiver $P_n$ we have
$\NS^{\G_m}((\P^1)^n)=\Pic^{\G_m}((\P^1)^n)\cong\Z^{n+1}$, 
where $(\theta_1,\dots,\theta_n)\in\Z^n$ defines an element
$\O_{(\P^1)^n}(\theta_1,\ldots,\theta_n)\in\Pic((\P^1)^n)$ and
for each $\O_{(\P^1)^n}(\theta_1,\ldots,\theta_n)$ the set of $\G_m$-linearisations
is a principal homogeneous space under the character group $\Z$ of $\G_m$.
Below we identify $\Pic^{\G_m}((\P^1)^n)$ with $H(P_n)\cap\Z^{n+2}$, where
$(\eta_1,\eta_2,\theta_1,\ldots,\theta_n)\in H(P_n)\cap\Z^{n+2}$ corresponds 
to the line bundle $\O_{(\P^1)^n}(\theta_1,\ldots,\theta_n)$ with a certain
$\G_m$-linearisation, and the $\G_m$-linearisations on 
$\O_{(\P^1)^n}(\theta_1,\ldots,\theta_n)$ corresponding 
to $(\eta_1,\eta_2,\theta_1,\ldots,\theta_n)$ and 
$(\eta_1',\eta_2',\theta_1,\ldots,\theta_n)$ differ by the character
$(\eta,-\eta)=(\eta_1-\eta_1',\eta_2-\eta_2')$ of $(\G_m)^2/\G_m\cong\G^m$.

\begin{prop}\label{prop:P_n-GITP1}
Let $R^*(P_n)=R(P_n)\setminus\bigcup_i\{s_i=0\}\cong(\A^2\setminus\{0\})^n$.%\\

\begin{enumerate}[label=(\alph*)]
\item\label{prop2.11_a} 
%{\rm (a)} 
The morphism $R^*(P_n)\to(\P^1)^n$ induces an isomorphism of stacks
\begin{equationi}\label{eq:isom-P_n-GITP1}
[R^*(P_n)/(((\G_m)^2\times(\G_m)^n)/\G_m)]\;\to\;[(\P^1)^n/\G_m].
\end{equationi}

\item\label{prop2.11_b} 
%{\rm (b)} 
Let $(\eta_1,\eta_2,\theta_1,\ldots,\theta_n)\in H(P_n)$ be integral. 
Then the line bundle $\O_{R^*(P_n)}$ with
$((\G_m)^2\times(\G_m)^n)/\G_m$-linearisation
given by $(\eta_1,\eta_2,\theta_1,\ldots,\theta_n)$ descends to the line bundle 
$\O_{(\P^1)^n}(\theta_1,\ldots,\theta_n)$ with a certain $\G_m$-linearisation, and both 
define isomorphic line bundles on the quotient stacks under the identification 
in~\eqref{eq:isom-P_n-GITP1}.
\end{enumerate}
\end{prop}

\begin{coro}\label{cor:reprPn-P1n-GIT}\ \\*[-1.2em]

\begin{enumerate}[label=(\alph*)]
\item\label{coro2.12_a} 
%{\rm (a)} 
For $\theta=(\eta_1,\eta_2,\theta_1,\ldots,\theta_n)\in H(P_n)$ such that 
$\theta_i>0$ for all $i$ a free representation $V=(s_1,\ldots,s_n)$ of $P_n$ is 
$\theta$-(semi)stable if and only if $s_i\cap 0=\emptyset$ for all $i$ and its 
image under~\eqref{eq:isom-Q_n-GITP1} is a (semi)stable section of $(\P^1)^n$
with respect to $\O_{(\P^1)^n}(\theta_1,\ldots,\theta_n)$ with the corresponding 
$\G_m$-linearisation. %\\

\item\label{coro2.12_b} 
%{\rm (b)} 
The image of $\interior C(P_n)\subset H(P_n)$ under the isomorphism 
$\Pic^{\G_m}((\P^1)^n)_\Q\cong H(P_n)$ is the $\G_m$-ample cone in 
$\Pic^{\G_m}((\P^1)^n)_\Q$. The GIT classes in these open cones coincide.%\\

\item\label{coro2.12_c} 
%{\rm (c)} 
For $\theta\in H(P_n)\cong\Pic^{\G_m}((\P^1)^n)_\Q$ such that 
$\theta_i>0$ for all $i$ the isomorphism~\eqref{eq:isom-Q_n-GITP1} induces 
isomorphisms of the stacks of $\theta$-(semi)stable points, of their moduli spaces 
and their inverse systems.
\end{enumerate}
\end{coro}

In the following, when considering free representations $V=(s_1,\ldots,s_n)$
of $Q_n$ and $P_n$ such that $s_i\cap 0=\emptyset$ up to isomorphism,
we will often write them as tuples of sections of $(\P^1)^n$.



\subsection{The functor of inverse limits of quiver varieties for \texorpdfstring{$P_n$}{Pn} and \texorpdfstring{$Q_n$}{QN}}
\label{subsec:functor-limPnQn}

Let $Q\!=\!Q_n$ or $Q\!=\!P_n$. Because all chambers are connected via the 
relative interiors of GIT equivalence classes of codimension one, we can rewrite 
the contravariant functor on $S$-schemes~\eqref{eq:functorinvlim} in 
proposition~\ref{prop:functorinvlim} as
\begin{equation}\label{eq:functorinvlim-QnPn}
Y\mapsto
\left\{\left.
(\phi_{V^\theta})_\theta\in\!\!\!\!\!\prod_{\text{$\theta$ generic}}\!\!\!\!\!
\mathcal M^\theta(Q)(Y)
\;\right|\!
\begin{array}{l}
\forall\,\theta,\theta'\,\text{generic}\;\;
\forall\:C_\tau\subset\interior C(Q)\;\text{GIT}\\
\text{equiv.\ class of codimension one such that}\\
\overline{C}_\tau\subseteq\overline{C}_\theta\cap\overline{C}_{\theta'}\colon
\;\phi_{\theta,\tau}\circ\phi_{V^{\theta}}=\phi_{\theta',\tau}\circ\phi_{V^{\theta'}}
\end{array}\!
\right\}
\end{equation}
\looseness-1
where the product is over representatives $\theta$ of the generic GIT equivalence classes
and $\phi_{V^\theta}\colon Y\to\mathcal M^\theta(Q)$ is the morphism determined by a
$\theta$-stable representation $V^\theta$ over $Y$.

%\medskip

For a representation $V$ of $Q$ over an $S$-scheme $Y$ stable with respect 
to some generic $\theta$ the polytopes $\Theta(V(y))$ for the fibres $V(y)$ 
over the geometric points $y$ of $Y$ defined in subsection~\ref{subsec:weightspacePnQn} 
have the property that $\{y'\in Y\:|\:\Theta(V(y))\subseteq \Theta(V(y'))\}$ are 
the points of an open subscheme of $Y$. Therefore for two representations $V,V'$ 
we have an open subscheme
\[
U(V,V')\;=\;\big\{\,y\in Y \mathbin{\big|}\interior\big(\Theta(V(y))\cap\Theta(V'(y))\big)
\neq\emptyset\,\big\}\;\subseteq \;Y.
\]
If $V^\theta,V^{\theta'}\!$ are part of a family over $Y$ satisfying the conditions
in~\eqref{eq:functorinvlim-QnPn}, then 
\[
U(V^\theta,V^{\theta'})
\;=\;\big\{\,y\in Y\mathbin{\big|}\text{$V^\theta(y)$ $\theta'$-stable}\,\big\}
\;=\;\big\{\,y\in Y\mathbin{\big|}\text{$V^{\theta'}\!(y)$ $\theta$-stable}\,\big\}
\] 
and $V^\theta|_{U(V^\theta,V^{\theta'})}\cong V^{\theta'}|_{U(V^\theta,V^{\theta'})}$
by the conditions coming from the walls meeting the interior of the polytopes
$\Theta(V^\theta(y))$, $\Theta(V^{\theta'}\!(y))$.

%\medskip

To study the conditions coming from walls that separate polytopes $\Theta(V(y))$ 
we consider the schemes of representations of $Q$ for weights in a GIT equivalence
class $C_\tau$ of codimension one with adjacent chambers $C_\theta,C_{\theta'}$. 
For $Q=Q_n$, $C_\tau\cap\Delta(2,n)$ is contained in a wall of the form 
$W_{\{J,J^\complement\}}$ for some $J\subset\{1,\ldots,n\}$, $2\leq|J|\leq n-2$. 
For $Q=P_n$, $C_\tau\cap\Delta^1\times\Delta^{n-1}$ is contained in a 
wall of the form $W_J$ for some $J\subset\{1,\ldots,n\}$, $1\leq|J|\leq n-1$. 
In case $Q=Q_n$ let $i\in J^\complement$, $j\in J$, in case $Q=P_n$ let $i=0,j=\infty$.
Let $Z_\tau\subset R^\tau(Q)$ be the closed subscheme of strictly 
$\tau$-semistable points and $V^\tau$ its tautological representation
considered as tuple of sections $(s^\tau_1,\ldots,s^\tau_n)$ of $\P^1_{Z_\tau}\to Z_\tau$.
The geometric fibres of $Z_\tau$ over $S$ consist of the following three strata: 
the subset where both conditions $s^\tau_l=s^\tau_j\Leftrightarrow l\in J$ and 
$s^\tau_k=s^\tau_i\Leftrightarrow k\in J^\complement$ 
are satisfied, and the two subsets $Z_{\tau,J},Z_{\tau,J^\complement}$ where only 
the first (resp.\ the second) condition is satisfied. 
In case $Q=Q_n$ assume $\sum_{j\in J}\theta_j>1$, then 
$\sum_{i\in J^\complement}\theta'_i>1$ and $Z_{\tau,J}=Z_{\tau}\cap R^{\theta'}(Q)$, 
$Z_{\tau,J^\complement}=Z_{\tau}\cap R^\theta(Q)$.
The image of $Z_{\tau,J}$ forms a projective space 
$\P^{|J^\complement|-2}_S\subset\mathcal M^{\theta'}(Q)$, the image of 
$Z_{\tau,J^\complement}$ forms a projective space 
$\P^{|J|-2}_S\subset\mathcal M^\theta(Q)$ and the image of $Z_\tau$ forms 
a subscheme $Z^\tau\subset\mathcal M^\tau(Q)$ isomorphic to the base scheme $S$. 
The morphisms
$\mathcal M^\theta(Q)\rightarrow\mathcal M^\tau(Q)\leftarrow\mathcal M^{\theta'}(Q)$
restrict to morphisms $\P^{|J|-2}_S\rightarrow Z^\tau\leftarrow\P^{|J^\complement|-2}_S$
and are isomorphisms elsewhere.
The case $Q=P_n$ is similar with projective spaces
$\P^{|J|-1}_S\subset\mathcal M^\theta(Q)$, 
$\P^{|J^\complement|-1}_S\subset\mathcal M^{\theta'}(Q)$. 

%\medskip

Coordinate functions on $R^*(Q)$ are given by the tautological family on $R^*(Q)$, 
considered as sections 
$s=(s_1\!=\!(s_{1,0}\!:\!s_{1,1}),\ldots,s_n\!=\!(s_{n,0}\!:\!s_{n,1}))$ 
of $\P^1_{R^*(Q)}\to R^*(Q)$, as follows. 
In case $Q=Q_n$ for $i,j\in\{1,\ldots,n\}$ over the invariant open subscheme 
$\{y\:|\:s_i(y)\neq s_j(y)\}\subset R^*(Q)$ there is the section
$\left(\begin{smallmatrix} s_{i,1}&-s_{i,0}\\ -s_{j,1}&s_{j,0}\end{smallmatrix}\right)$
in $\PGL(2)$ that transforms the tautological family $s$ to the family $\tilde{s}$ 
with $\tilde{s}_i=(0:1),\tilde{s}_j=(1:0)$.
In case $Q=P_n$ we have the additional sections $s_i=s_0=(0:1)$, $s_j=s_\infty=(1:0)$ 
and we set $\tilde{s}=s$.
The invariant open subscheme 
\[
U_\tau=\big\{\,y\:\big|\:s_i(y)\neq s_j(y),\:
\forall\,k\!\in\!J^\complement\colon s_k(y)\neq s_j(y),\:
\forall\,l\!\in\!J\colon s_l(y)\neq s_i(y)\big\}\subset R^*(Q)
\]
contains $Z_\tau$. The algebra of invariant functions on $U_\tau$ corresponds to 
the subalgebra of torus invariant functions in the $\O_S$-algebra generated by 
$\frac{\tilde{s}_{l,1}}{\tilde{s}_{l,0}}$ for $l\in J$ and
$\frac{\tilde{s}_{k,0}}{\tilde{s}_{k,1}}$ for $k\in J^\complement$. 
This $\O_S$-algebra of torus invariants is generated by 
$f^{i,j}_{k,l}=\frac{\tilde{s}_{k,0}\tilde{s}_{l,1}}{\tilde{s}_{k,1}\tilde{s}_{l,0}}$
for $l\in J$, $k\in J^\complement$.
The function $f^{i,j}_{k,l}$ can be written in terms of $s$ as
\[
f^{i,j}_{k,l}=\frac{(s_{j,0}s_{l,1}-s_{j,1}s_{l,0})(s_{i,1}s_{k,0}-s_{i,0}s_{k,1})}
{(s_{i,1}s_{l,0}-s_{i,0}s_{l,1})(s_{j,0}s_{k,1}-s_{j,1}s_{k,0})}.
\]
These invariant functions are sometimes called the cross-ratios of four sections.
Since invariant regular functions on the representation space correspond to regular 
functions on the moduli spaces we obtain the following result.

\begin{lemma}\ \\*[-1.2em]\label{le:fijlk}

\begin{enumerate}[label=(\alph*)]
\item\label{lemma2.13_a}
%{\rm(a)} 
The invariant regular functions $f^{i,j}_{k,l}$ on $U_\tau$ define 
regular functions on the image of $U_\tau\cap R^\tau(Q)$ (resp.\ $U_\tau\cap R^\theta(Q)$,
$U_\tau\cap R^{\theta'}(Q)$) in $\mathcal M^\tau(Q)$ (resp.\ $\mathcal M^\theta(Q)$, 
$\mathcal M^{\theta'}(Q)$) which we also denote $f^{i,j}_{k,l}$.
These satisfy $\phi_{\theta,\tau}^*(f^{i,j}_{k,l})=f^{i,j}_{k,l}$,
$\phi_{\theta',\tau}^*(f^{i,j}_{k,l})=f^{i,j}_{k,l}$. %\\

\item\label{lemma2.13_b}
%{\rm(b)} 
For $z\in Z^\tau$ the local ring $\O_{\mathcal M^\tau(Q),z}$ is the 
localisation of an $\O_S$-algebra generated by the functions $f^{i,j}_{k,l}$ for 
$l\in J$, $k\in J^\complement$.
\end{enumerate}
\end{lemma}

We express the conditions in~\eqref{eq:functorinvlim-QnPn} for a family 
$(V^\theta)_\theta$ locally in terms of equations.

\begin{lemma}\label{le:eq-invlimit}
Let $V^\theta,V^{\theta'}$ be free representations of $Q=Q_n$ or $Q=P_n$ over an %\linebreak 
\text{$S$-scheme} $Y$ stable with respect to generic $\theta,\theta'$. 
Assume $V^\theta,V^{\theta'}$ are part of a collection as in 
\eqref{eq:functorinvlim-QnPn}. 
We consider $V^\theta,V^{\theta'}$ as tuples of sections 
$(s^\theta_i)_i=(s^\theta_{i,0}:s^\theta_{i,1})_i$, 
$(s^{\theta'}_i)_i=(s^{\theta'}_{i,0}:s^{\theta'}_{i,1})_i$ of $\P^1_Y$, 
in case $Q=P_n$ we add the sections $s^\theta_0,s^{\theta'}_0=(0:1)$,
$s^\theta_\infty,s^{\theta'}_\infty=(1:0)$.
In case $Q=Q_n$ let $i,j\in\{1,\ldots,n\}$, in case $Q=P_n$ let $i=0,j=\infty$.
Let $U_{i,j}=\{y\:|\:s^\theta_i(y)\neq s^\theta_j(y),
s^{\theta'}_i(y)\neq s^{\theta'}_j(y)\}\subseteq Y$.
We choose homogeneous coordinates $x^\theta_0,x^\theta_1$ and 
$x^{\theta'}_0,x^{\theta'}_1$ of $\P^1_{U_{i,j}}$ 
such that $s^\theta_i=(0:1),s^\theta_j=(1:0)$ with respect to $x^\theta_0,x^\theta_1$
and $s^{\theta'}_i=(0:1), s^{\theta'}_j=(1:0)$ with respect to $x^{\theta'}_0,x^{\theta'}_1$. 
Then $V^\theta$ and $V^{\theta'}$ are over $U_{i,j}$ related by the equations
\begin{equation}\label{eq:theta-theta'}
s^\theta_{k,0}s^\theta_{l,1}s^{\theta'}_{k,1}s^{\theta'}_{l,0}
=s^\theta_{k,1}s^\theta_{l,0}s^{\theta'}_{k,0}s^{\theta'}_{l,1}
\end{equation}
for all $k,l\in\{1,\ldots,n\}$.
\end{lemma}
\begin{proof}
Over $U(V^\theta,V^{\theta'})$ the equations hold because
$V^\theta|_{U(V^\theta,V^{\theta'})}\cong V^{\theta'}|_{U(V^\theta,V^{\theta'})}$.

Consider the local situation around a point 
$y\in U_{i,j}\setminus U(V^\theta,V^{\theta'})$. 
Assume first that the polytopes $\Theta(V^\theta(y))$ and $\Theta(V^{\theta'}(y))$ 
meet in an inner wall $W$. 
In case %\linebreak 
\text{$Q=P_n$} let $W=W_J$ and we can assume that 
$s^\theta_h(y)=s^\theta_j(y)=s^\theta_\infty(y)=(1:0)$ for $h\in J$ and 
$s^{\theta'}_h(y)=s^{\theta'}_i(y)=s^{\theta'}_0(y)=(0:1)$ for $h\in J^\complement$. 
In case $Q=Q_n$ let $W=W_{\{J,J^\complement\}}$ and we can assume 
$j\in J$, $i\in J^\complement$ and $s^\theta_h(y)=s^\theta_j(y)=(1:0)$ 
for $h\in J$, $s^{\theta'}_h(y)=s^{\theta'}_i(y)= (0:\nobreak 1)$ for $h\in J^\complement$.
For $k,l$ there are the cases: $k,l\in J$ (similar: $k,l\in J^\complement$)
or $k\in J,l\in J^\complement$ (similar: $k\in J^\complement,l\in J$).

\emph{Case $k\in J^\complement,l\in J$}:
We have $s^\theta_{l,0}(y),s^{\theta'}_{l,0}(y)\neq0$ and 
$s^\theta_{k,1}(y),s^{\theta'}_{k,1}(y)\neq0$.
In a neighbourhood of $y$ equation~\eqref{eq:theta-theta'} is equivalent to the 
equation $f^{\theta,i,j}_{k,l}=f^{\theta',i,j}_{k,l}$ with 
$f^{\theta,i,j}_{k,l}=\frac{s^\theta_{k,0}s^\theta_{l,1}}{s^\theta_{k,1}s^\theta_{l,0}}$,
$f^{\theta',i,j}_{k,l}=\frac{s^{\theta'}_{k,0}s^{\theta'}_{l,1}}
{s^{\theta'}_{k,1}s^{\theta'}_{l,0}}$.
The regular functions $f^{\theta,i,j}_{k,l}$ (resp.\ $f^{\theta',i,j}_{k,l}$) 
are pullbacks of the regular functions $f^{i,j}_{k,l}$ on
$\mathcal M^\theta(Q)$ (resp.\ $\mathcal M^{\theta'}(Q)$) via 
$\phi_{V^\theta}$ (resp.\ $\phi_{V^{\theta'}}$). 
Because $\phi_{\theta,\tau}\circ\phi_{V^\theta}=\phi_{\theta',\tau}\circ\phi_{V^{\theta'}}$
and using lemma~\ref{le:fijlk}.\ref{lemma2.13_a} it follows $f^{\theta,i,j}_{k,l}=f^{\theta',i,j}_{k,l}$.

\emph{Case $k,l\in J$}:
We have $s^\theta_{k,0}(y),s^{\theta'}_{k,0}(y)\neq0$ and 
$s^\theta_{l,0}(y),s^{\theta'}_{l,0}(y)\neq0$.
In a neighbourhood of $y$ equation~\eqref{eq:theta-theta'} is equivalent to the equation
$\frac{s^\theta_{l,1}}{s^\theta_{l,0}}\frac{s^{\theta'}_{k,1}}{s^{\theta'}_{k,0}}
=\frac{s^\theta_{k,1}}{s^\theta_{k,0}}\frac{s^{\theta'}_{l,1}}{s^{\theta'}_{l,0}}$.
Multiplying with $\frac{s^\theta_{g,0}s^{\theta'}_{h,0}}{s^\theta_{g,1}s^{\theta'}_{h,1}}$
for $g\in J^\complement$ such that $s^\theta_g(y)\neq(0:1)$ and $h\in J$ such that 
$s^{\theta'}_h(y)\neq(1:0)$, we obtain the equivalent equation 
$f^{\theta,i,j}_{g,l}f^{\theta',i,j}_{h,k}=f^{\theta,i,j}_{g,k}f^{\theta',i,j}_{h,l}$.
This equation holds because by the first case we have 
$f^{\theta,i,j}_{g,l}=f^{\theta',i,j}_{g,l}$, $f^{\theta,i,j}_{g,k}=f^{\theta',i,j}_{g,k}$.

This shows that equations~\eqref{eq:theta-theta'} hold for $V^\theta,V^{\theta'}$
in a neighbourhood of $y$ if the polytopes $\Theta(V^\theta(y))$ and 
$\Theta(V^{\theta'}(y))$ have overlapping interiors or meet in an inner wall.

To show the general case we show that if the union of the sets of walls separating
$\Theta(V^\theta(y))$, $\Theta(V^{\theta'}(y))$ and 
$\Theta(V^{\theta'}(y))$, $\Theta(V^{\theta''}(y))$ are exactly the walls that separate 
$\Theta(V^\theta(y)),\Theta(V^{\theta''}(y))$ and if $V^\theta,V^{\theta'}$ and
$V^{\theta'},V^{\theta''}$ in a neighbourhood of $y$ are related by equations of 
the form~\eqref{eq:theta-theta'}, then so are $V^\theta,V^{\theta''}$.
Assume that $\Theta(V^\theta(y))$ and $\Theta(V^{\theta'}(y))$ are separated by an
inner wall $W_J$ (resp.\ $W_{\{J,J^\complement\}}$) and that $\Theta(V^{\theta'}(y))$ 
and $\Theta(V^{\theta''}(y))$ are separated by an inner wall $W_{J'}$ (resp.\ 
$W_{\{J',(J')^\complement\}}$) such that $J\subset J'$, these walls are part 
of the boundary of $\Theta(V^{\theta'}(y))$ and both walls separate 
$\Theta(V^\theta(y))$ and $\Theta(V^{\theta''}(y))$. 
If $s_i^\theta(y)\neq s_j^\theta(y)$, $s_i^{\theta''}(y)\neq s_j^{\theta''}(y)$
then either $i\in J$, $j\in(J')^\complement$ or $j\in J$, $i\in(J')^\complement$ 
and thus $s_i^{\theta'}(y)\neq s_j^{\theta'}(y)$. 
We assume $j\in J$, $i\in(J')^\complement$.

\emph{Case $k\in J',l\in J^\complement$ (similar: $k\in (J')^\complement,l\in J$)}.
It is $s^{\theta'}_k(y)\neq (0:1)$, $s^{\theta'}_l(y)\neq(1:0)$ and in 
a neighbourhood of $y$ we have 
\[
s^\theta_{k,0}s^\theta_{l,1}
s^{\theta''}_{k,1}s^{\theta''}_{l,0}s^{\theta'}_{k,0}s^{\theta'}_{l,1}
=s^\theta_{k,0}s^\theta_{l,1}
s^{\theta'}_{k,1}s^{\theta'}_{l,0}s^{\theta''}_{k,0}s^{\theta''}_{l,1}
=s^{\theta'}_{k,0}s^{\theta'}_{l,1}s^\theta_{k,1}s^\theta_{l,0}
s^{\theta''}_{k,0}s^{\theta''}_{l,1},
\]
thus 
$s^\theta_{k,0}s^\theta_{l,1}s^{\theta''}_{k,1}s^{\theta''}_{l,0}
=s^\theta_{k,1}s^\theta_{l,0}s^{\theta''}_{k,0}s^{\theta''}_{l,1}$.

\emph{Case $k\in J',l\in J$ (similar: $k\in(J')^\complement, l\in J^\complement$)}.
Choose $h\in J^\complement\cap J'$.
Then $s^{\theta'}_k(y),s^{\theta'}_l(y)\neq (0:1)$ and
$s^\theta_h(y)=(0:1),s^{\theta''}_h(y)=(1:0)$, 
$s^{\theta'}_h(y)\neq (1:0),(0:1)$.
In a neighbourhood of $y$ we have
\[
\begin{aligned}
s^\theta_{k,0}s^\theta_{l,1}s^{\theta''}_{k,1}s^{\theta''}_{l,0}
s^{\theta'}_{k,0}s^{\theta'}_{l,0}
s^{\theta'}_{h,1}s^\theta_{h,1}s^{\theta''}_{h,0}
&=s^\theta_{k,0}s^\theta_{l,1}s^{\theta'}_{k,1}s^{\theta''}_{l,0}
s^{\theta''}_{k,0}s^{\theta'}_{l,0}
s^{\theta''}_{h,1}s^\theta_{h,1}s^{\theta'}_{h,0}\\
&=s^{\theta'}_{k,0}s^\theta_{l,1}s^\theta_{k,1}s^{\theta''}_{l,0}
s^{\theta''}_{k,0}s^{\theta'}_{l,0}
s^{\theta''}_{h,1}s^{\theta'}_{h,1}s^\theta_{h,0}\\
&=s^{\theta'}_{k,0}s^{\theta'}_{l,1}s^\theta_{k,1}s^{\theta''}_{l,0}
s^{\theta''}_{k,0}s^\theta_{l,0}
s^{\theta''}_{h,1}s^\theta_{h,1}s^{\theta'}_{h,0}\\
&=s^{\theta'}_{k,0}s^{\theta''}_{l,1}s^\theta_{k,1}s^{\theta'}_{l,0}
s^{\theta''}_{k,0}s^\theta_{l,0}
s^{\theta'}_{h,1}s^\theta_{h,1}s^{\theta''}_{h,0}
\end{aligned}
\]
thus 
$s^\theta_{k,0}s^\theta_{l,1}s^{\theta''}_{k,1}s^{\theta''}_{l,0}
=s^\theta_{k,1}s^\theta_{l,0}s^{\theta''}_{k,0}s^{\theta''}_{l,1}$.
\end{proof}


\begin{prop}\label{prop:functorinvlimQnPn}
For $Q=Q_n$ or $Q=P_n$ the functor of the inverse limit of moduli spaces
of representations is isomorphic to the functor $\Lim(Q)$ defined by
\begin{equation}\label{eq:functorinvlimQnPn}
\;Y\;\mapsto\;
\Bigg\{
(\phi_{V^\theta})_\theta\,\in\!\!\!\prod_{\text{$\theta$ generic}}\!\!\!\!
\mathcal M^\theta(Q)(Y)
\;\Bigg|\;
\text{locally equations~\eqref{eq:theta-theta'} hold for $(V^\theta)_\theta$}\:
\Bigg\}
\end{equation}
where the product is over representatives $\theta$ of the generic GIT equivalence classes
and $\phi_{V^\theta}\colon Y\to\mathcal M^\theta(Q)$ is the morphism determined by a
$\theta$-stable representation $V^\theta$ over $Y$.
\end{prop}

\begin{proof}
By lemma~\ref{le:eq-invlimit} for a family $(V^\theta)_\theta$ 
satisfying the conditions of~\eqref{eq:functorinvlim-QnPn} the equations 
\eqref{eq:theta-theta'} hold locally. We show the opposite implication.

Let $(V^\theta)_\theta$ be a family of representations over an $S$-scheme $Y$
as in~\eqref{eq:functorinvlimQnPn}. Let $y\in Y$. 
Let $C_\tau$ be an equivalence class of codimension $1$ in the interior 
and $\theta,\theta'$ generic such that
$\overline{C}_\tau=\overline{C}_\theta\cap\overline{C}_{\theta'}$.
We can consider $V^\theta,V^{\theta'}$ in a neighbourhood $U$ of $y$ as tuples of 
sections $(s^\theta_i)_i,(s^{\theta'}_i)_i$ of $\P^1_U\to U$.
In case $Q=Q_n$ let $i,j\in\{1,\ldots,n\}$ such that $s^\theta_i(y)\neq s^\theta_j(y)$ 
and $s^{\theta'}_i(y)\neq s^{\theta'}_j(y)$ and choose coordinates such that 
$s^\theta_i,s^{\theta'}_i=(0:1)$ and $s^\theta_j,s^{\theta'}_j=(1:0)$. 
In case $Q=P_n$ let $i=0,j=\infty$, then $s^\theta_i=s^\theta_0=(0:1)$, %\linebreak
$s^\theta_j=s^\theta_\infty=(1:0)$ and the same for $\theta'$.

If there exists $k\in\{1,\ldots,n\}$ such that $s^\theta_k(y)\neq(0:1),(1:0)$,
$s^{\theta'}_k(y)\neq(0:1),(1:0)$ then the equations 
$s^\theta_{k,0}s^\theta_{l,1}s^{\theta'}_{k,1}s^{\theta'}_{l,0}
=s^\theta_{k,1}s^\theta_{l,0}s^{\theta'}_{k,0}s^{\theta'}_{l,1}$
show that in a neighbourhood of $y$ $V^\theta\cong V^{\theta'}$ and thus
$\phi_{\theta,\tau}\circ\phi_{V^\theta}=\phi_{\theta',\tau}\circ\phi_{V^{\theta'}}$.

Otherwise there is $J\subset\{1,\ldots,n\}$, where we can assume that $j\in J$,
$i\in J^\complement$, such that 
$\forall\,k\in J^\complement\colon s^\theta_k(y)=(0:1),\:s^{\theta'}_k(y)\neq(1:0)$ 
and $\forall\,l\in J\colon s^\theta_l(y)\neq (0:1)$, $s^{\theta'}_l(y)=(1:0)$.
Let $x\in S$ be the image of $y\in Y$ and $z\in Z^\tau$ the point over $x$, then 
$\phi_{\theta,\tau}\circ\phi_{V^\theta}(y)=z=\phi_{\theta',\tau}
\circ\phi_{V^{\theta'}}(y)$.
As equations~\eqref{eq:theta-theta'} hold, the functions
$f^{\theta,i,j}_{k,l}=\frac{s^\theta_{k,0}s^\theta_{l,1}}{s^\theta_{k,1}s^\theta_{l,0}}$ 
and $f^{\theta'\!,i,j}_{k,l}=\frac{s^{\theta'}_{k,0}s^{\theta'}_{l,1}}
{s^{\theta'}_{k,1}s^{\theta'}_{l,0}}$ coincide. We may consider these functions 
as elements of the local ring $\O_{Y,y}$.
Using lemma~\ref{le:fijlk}, these are the pullbacks of the functions $f^{i,j}_{k,l}$ 
in $\O_{\mathcal M^\tau(Q),z}$ under $\phi_{\theta,\tau}\circ\phi_{V^\theta}$ and
$\phi_{\theta',\tau}\circ\phi_{V^{\theta'}}$, and it follows that the homomorphisms 
of local rings $\O_{\mathcal M^\tau(Q),z}\to\O_{Y,y}$ agree.

Since $\phi_{\theta,\tau}\circ\phi_{V^\theta}$ and 
$\phi_{\theta',\tau}\circ\phi_{V^{\theta'}}$ coincide as maps of sets of points and 
in each point the induced homomorphisms of local rings coincide, it follows
$\phi_{\theta,\tau}\circ\phi_{V^\theta}=\phi_{\theta',\tau}\circ\phi_{V^{\theta'}}$.
\end{proof}

\begin{thebibliography}{99}
\bibitem[1]{Al14} Jarod Alper, ``Adequate moduli spaces and geometrically reductive group schemes'', \emph{Algebr. Geom.}\textbf{1}(2014),no.4,489--531, \url{https://arxiv.org/abs/1005.2398}.
\bibitem[2]{An73} Sivaramakrishna Anantharaman, ``Schémas en groupes, espaces homogènes et espaces algébriques sur une base de dimension1'', in \emph{Sur les groupes algébriques}, Bull. Soc. Math. France,Mém.33, Soc. Math. France,1973,5--79.
\bibitem[4]{Br14} Sylvain Brochard, ``Topologies de Grothendieck, descente, quotients'', in \emph{Autour des schémas en groupes.Vol.I}, Panor.Synthèses42/43, Soc. Math. France,Paris,2014,1--62, \url{https://arxiv.org/abs/1210.0431}.
\bibitem[5]{Ch08} Calin Chindris, ``Notes on GIT-fans for quivers'', \url{https://arxiv.org/abs/0805.1440}.
\bibitem[6]{Cr96} William Crawley-Boevey, ``Subrepresentations of general representations of quivers'', \emph{Bull. London Math. Soc.}\textbf{28}(1996),no.4,363--366.
\bibitem[7]{DH98} Igor V. Dolgachev and Yi Hu, ``Variation of geometric invariant theory quotients'', \emph{Publ. Math.,Inst. Hautes Étud.Sci.}(1998),no.87,5--51, \url{https://arxiv.org/abs/alg-geom/9402008}.
\bibitem[8]{GGMS87} Israel M. Gel'fand, Mark Goresky, Robert D. MacPherson and Vera V. Serganova, ``Combinatorial geometries, convex polyhedra, and Schubert cells'', \emph{Adv. in Math.}\textbf{63}(1987),no.3,301--316.
\bibitem[12]{EGA} Alexander Grothendieck and Jean A. Dieudonné, \emph{Éléments de géométrie algébrique.I--IV}, Publ. Math.,Inst. Hautes Étud.Sci.4,8,11,17,20,24,28,32(1960--1967).
\bibitem[13]{EGA1} Alexander Grothendieck and Jean A. Dieudonné, \emph{Éléments de géométrie algébrique.I}, Grundlehren der Mathematischen Wissenschaften166, Springer-Verlag,Berlin,1971.
\bibitem[16]{HP02} Lutz Hille and José Antonio de la Peña, ``Stable representations of quivers'', \emph{J. Pure Appl. Algebra}\textbf{172}(2002),no.2-3,205--224.
\bibitem[17]{Ka93a} Mikhail M. Kapranov, ``Chow quotients of Grassmannians.I'', in \emph{I.M.Gel'fand Seminar}, Adv. Soviet Math.16, Amer. Math. Soc.,Providence,RI,1993,29--110, \url{https://arxiv.org/abs/alg-geom/9210002}.
\bibitem[20]{Ki94} Alastair D. King, ``Moduli of representations of finite-dimensional algebras'', \emph{Quart. J. Math. Oxford Ser.}(2)\textbf{45}(1994),no.180,515--530.
\bibitem[22]{La69} Daniel Lazard, ``Autour de la platitude'', \emph{Bull. Soc. Math. France}\textbf{97}(1969),81--128.
\bibitem[24]{MFK} David Mumford, John Fogarty and Frances Kirwan, \emph{Geometric invariant theory}, third ed., Ergebnisse der Mathematik und ihrer Grenzgebiete(2)34, Springer-Verlag,Berlin,1994.
\bibitem[25]{Re00} Nicolas Ressayre, ``The GIT-equivalence for G-line bundles'', \emph{Geom. Dedicata}\textbf{81}(2000),no.1-3,295--324, \url{https://arxiv.org/abs/math/9811053}.
\bibitem[26]{Scho92} Aidan Schofield, ``General representations of quivers'', \emph{Proc. London Math. Soc.}(3)\textbf{65}(1992),no.1,46--64.
\bibitem[29]{Ses77} Conjeeveram S. Seshadri, ``Geometric reductivity over arbitrary base'', \emph{Advances in Math.}\textbf{26}(1977),no.3,225--274.
\bibitem[30]{SP} Stacks Project Authors, \emph{Stacks project}, \url{https://stacks.math.columbia.edu}.
\bibitem[31]{Th96} Michael Thaddeus, ``Geometric invariant theory and flips'', \emph{J. Amer. Math. Soc.}\textbf{9}(1996),no.3,691--723, \url{https://arxiv.org/abs/alg-geom/9405004}.
\end{thebibliography}
\end{document}
